\documentclass[11pt,letterpaper]{article}
\usepackage[T1]{fontenc}
\usepackage[utf8]{inputenc}
\usepackage{lmodern}
\usepackage[margin=1in,headheight=14pt]{geometry}
\usepackage{amsmath,amssymb,amsthm,mathtools}
\usepackage{microtype,booktabs,longtable,array,enumitem}
\usepackage[dvipsnames]{xcolor}
\usepackage{fancyhdr}
\usepackage{hyperref}
\usepackage[nameinlink,noabbrev]{cleveref}
\hypersetup{colorlinks=true,linkcolor=MidnightBlue,citecolor=MidnightBlue,urlcolor=MidnightBlue,
 pdftitle={Surquaternions: Algebra, Infinitesimal Geometry, and Support-Controlled Analysis},
 pdfauthor={OpenAI, prepared for Vladimir Reshetnikov}}
\pagestyle{fancy}\fancyhf{}
\fancyhead[L]{\small Surquaternions}
\fancyhead[R]{\small Algebra, geometry, and analysis}
\fancyfoot[C]{\thepage}
\setlength{\parindent}{1.2em}
\setlength{\parskip}{3pt}
\setlist{nosep,leftmargin=*}
\setcounter{tocdepth}{2}
\newtheorem{theorem}{Theorem}[section]
\newtheorem{proposition}[theorem]{Proposition}
\newtheorem{lemma}[theorem]{Lemma}
\newtheorem{corollary}[theorem]{Corollary}
\theoremstyle{definition}
\newtheorem{definition}[theorem]{Definition}
\newtheorem{example}[theorem]{Example}
\newtheorem{remark}[theorem]{Remark}
\newtheorem{warning}[theorem]{Warning}
\newcommand{\No}{\mathbf{No}}
\newcommand{\R}{\mathbb R}
\newcommand{\C}{\mathbb C}
\newcommand{\Q}{\mathbb Q}
\newcommand{\Z}{\mathbb Z}
\newcommand{\N}{\mathbb N}
\newcommand{\HH}{\mathbb H}
\newcommand{\SQ}{\mathbb H_{\No}}
\newcommand{\HF}{\mathbb H_F}
\newcommand{\DG}{D_\Gamma}
\newcommand{\KG}{K_\Gamma}
\newcommand{\OO}{\mathcal O}
\newcommand{\mm}{\mathfrak m}
\newcommand{\ii}{\mathbf i}
\newcommand{\jj}{\mathbf j}
\newcommand{\kk}{\mathbf k}
\newcommand{\Ss}{\mathbb S_F}
\newcommand{\val}{\nu}
\newcommand{\abs}[1]{\left|#1\right|}
\newcommand{\norm}[1]{\left\lVert#1\right\rVert}
\newcommand{\br}[2]{[#1,#2]}
\DeclareMathOperator{\Rea}{Re}
\DeclareMathOperator{\Imag}{Im}
\DeclareMathOperator{\st}{st}
\DeclareMathOperator{\supp}{supp}
\DeclareMathOperator{\Ad}{Ad}
\DeclareMathOperator{\ad}{ad}
\DeclareMathOperator{\Sp}{Sp}
\DeclareMathOperator{\SO}{SO}
\DeclareMathOperator{\SU}{SU}
\DeclareMathOperator{\GL}{GL}
\DeclareMathOperator{\fin}{fin}
\DeclareMathOperator{\Sin}{Sin}
\DeclareMathOperator{\Cos}{Cos}
\DeclareMathOperator{\Oz}{Oz}
\DeclareMathOperator{\BCH}{BCH}
\newcommand{\ExpC}{\operatorname{Exp}_{\mathrm{can}}}
\newcommand{\exph}{\operatorname{exp}_{\mathrm H}}
\newcommand{\logh}{\operatorname{log}_{\mathrm H}}
\newcommand{\smallo}{\mathrm{o}}
\newcommand{\bigO}{\mathrm{O}}
\newcommand{\repo}{\texttt{VladimirReshetnikov/Surreal}}
\title{\textbf{Surquaternions}\\[5pt]
\Large Algebra, Infinitesimal Geometry, and\\Support-Controlled Analysis over the Surreal Numbers}
\author{A mathematical development prepared for Vladimir Reshetnikov\\
\small OpenAI}
\date{21 September 2026}
\begin{document}
\maketitle
\begin{abstract}
A surquaternion is an expression $a+b\ii+c\jj+d\kk$ with surreal coefficients
and Hamilton's multiplication rules. This elementary definition leads to a
substantial theory, provided three structures are kept separate: algebra over a
real closed field, the surreal normal-form valuation, and analytic operations
whose infinite sums have certified supports. We develop all three.

The algebraic part establishes division, conjugacy and complex slices,
rotations, a constructive fundamental theorem for one-sided polynomials,
and finite-dimensional spectral theory. The valuation identifies surquaternions
with set-supported Hahn series having ordinary quaternion coefficients. It gives
an exact residue algebra, a multiplicative leading term, a support-controlled
exponential and logarithm, and a filtered nonabelian group of infinitesimal
rotations. A multivariable lifting theorem supplies quaternionic implicit
solutions without assuming convergence of Newton sequences. We calculate the
conditioning of square-root lifting explicitly.

For analysis, we distinguish slice regularity, Fueter regularity, and derivations
of the scalar field. We construct local slice series, prove their representation
and product formulas, and give coefficientwise Cauchy and Fueter-transfer
results. Finally, the established Ehrlich--Kaplan scalar phase normalization
produces a global conjugation-equivariant quaternionic exponential. Its complete
logarithm fibres and its differential are computed, revealing critical spheres
at infinite surreal radii. Throughout, full-class fine topology is distinguished
from the intrinsic topology of a set-sized Hahn field. Classical results are
identified as such; the article claims a coherent development, not a priority
claim for a new number system or a solution of a named open conjecture.
\end{abstract}
\vspace{6pt}
\noindent\textbf{Keywords.} Surreal numbers; quaternion division algebras; Hahn
series; slice regularity; infinitesimal rotations; noncommutative lifting.
\newpage
\tableofcontents
\newpage

\section{Scope, provenance, and the main distinctions}
\label{sec:scope}
The purpose of this article is not merely to replace each real coefficient of a
quaternion by a surreal number. That replacement settles the multiplication
rule, but leaves the interesting questions unanswered. Which polynomial
identities survive? What is the leading term of a noncommutative product? Can a
rotation carry several incomparable infinitesimal scales? Which exponential is
meant at an infinite imaginary argument? What, precisely, does an analytic
infinite sum denote?

This article follows the architecture of \repo.
The repository was inspected at commit
\texttt{39f2be6667ade51bca2b45daa47e289d69c09764}. Its README and document catalogue
separate finite algebra, Hahn summation, and common-domain analytic constructions;
the present development continues that separation \cite{repo}. The inspected
README explicitly says that its generic Lean prerequisites do not yet construct
the surreal field or its normal-form bridge. We therefore do not treat the
repository's build as a formal verification of the theorems below.

\subsection{What is classical and what is developed here}
Hamilton's algebra over an ordered field is standard quaternion algebra
\cite{voight}. Its specialization to $\No$ does not make the division formula,
conjugacy classification, or spectral theorem new. Likewise, the polynomial
star product, spherical zero sets, and stem-function construction belong to the
established theory of quaternionic slice regularity \cite{gp}. The scalar global
phase normalization used in \cref{sec:global} is due to Ehrlich and Kaplan
\cite{ek}, not introduced here.

The article's additional work is to fit these structures together with exact
surreal support hypotheses, prove the resulting statements, and make the
non-Archimedean and proper-class qualifications explicit. Particularly useful
combined results are the support-localization theorem, the filtered rotation
calculus, support-certified implicit lifting, the square-root conditioning
formula, and the differential and fibres of the radial global exponential.
These are presented as results proved in this article, not as exhaustively
verified claims of first publication. Literature searches were targeted, not a
complete bibliographic classification of every use of quaternion algebras over
non-Archimedean fields.

\subsection{Three levels of scalars}
Most finite algebra is proved over a real closed ordered field $F$. This is the
right level of generality: the proof then applies both to $\No$ and to suitable
set-sized subfields. Analytic constructions requiring infinite sums are placed
in a specified Hahn workspace
\[
 K_\Gamma=\R((t^\Gamma)),\qquad D_\Gamma=\HH(K_\Gamma),
\]
where $\Gamma$ is a set-sized divisible ordered subgroup of the additive surreal
group and $t^\gamma=\omega^{-\gamma}$. A statement involving all surreals is read
as a class statement, or as a scheme in a set theory with definable classes.

\begin{center}
\begin{tabular}{@{}>{\raggedright\arraybackslash}p{.22\linewidth}>{\raggedright\arraybackslash}p{.33\linewidth}>{\raggedright\arraybackslash}p{.36\linewidth}@{}}
\toprule
Level & Available structure & What must not be inferred \\
\midrule
$\HH(F)$ & Finite algebra, order-valued norm, polynomial geometry & Real compactness or unrestricted analytic transfer \\
$D_\Gamma$ & Set-sized supports, valuation, Hahn evaluation & That every power series is evaluable everywhere \\
$\SQ$ & All surreal scales and normal forms & That ordinary convergent sequences model fine limits \\
\bottomrule
\end{tabular}
\end{center}

\subsection{Reading guide}
Sections 2--6 contain the algebraic theory. Sections 7--9 explain valuation and
infinitesimal Lie theory. Sections 10--12 construct and analyze global phases.
Sections 13--15 develop local analytic operations, implicit lifting, and scalar
derivations. The final sections give implementation guidance and delimit the
claims. The appendices collect notation and verification information.

\section{Foundations and the definition}
\label{sec:foundations}
We use two established facts about Conway's surreal numbers: $\No$ is a real
closed ordered class field, and every surreal has a unique normal form with
real coefficients and set-sized reverse-well-ordered exponents
\cite{conway,gonshor,ek}. No proof here uses an unrestricted sum indexed by a
proper class.

\begin{definition}[Surquaternions]
For an ordered field $F$, let
\[
 \HH(F)=F\oplus F\ii\oplus F\jj\oplus F\kk,
 \qquad \ii^2=\jj^2=\kk^2=\ii\jj\kk=-1.
\]
The \emph{surquaternions} are $\SQ=\HH(\No)$. We write $\HH=\HH(\R)$ for the
ordinary quaternions. The four coefficients are central: they commute with the
basis elements and with every surquaternion.
\end{definition}
Thus a surquaternion is an ordered quadruple of surreals, not a new Conway cut
with quaternionic left and right options. The order used to construct surreal
numbers is an order on their scalar components. It does not become a compatible
ordering of the quaternion division algebra.

For $q=a+b\ii+c\jj+d\kk$, put
\[
 \bar q=a-b\ii-c\jj-d\kk,\quad
 \Rea q=a,\quad \Imag q=b\ii+c\jj+d\kk,\quad
 N(q)=a^2+b^2+c^2+d^2.
\]
The notation $\Rea q$ means the \emph{central scalar part}; it need not be a real
number. The quaternion conjugation is distinct from a derivation, a standard
part map, or conjugation in a chosen complex slice.

\begin{proposition}[Associativity and change of scalars]
$\HH(F)$ is an associative unital $F$-algebra of dimension four. Every ordered
field embedding $F\hookrightarrow E$ induces an injective algebra homomorphism
$\HH(F)\hookrightarrow\HH(E)$ preserving conjugation and norm square.
\end{proposition}
\begin{proof}
One may impose Hamilton's relations in the free associative algebra and reduce
every word to the span of $1,\ii,\jj,\kk$. Independence, and hence the asserted
multiplication, follows from the faithful matrix representation in
\cref{prop:matrix}. Alternatively, the coordinate multiplication below verifies
associativity directly as a polynomial identity. Applying the scalar embedding
to each of four coordinates proves the last assertion.
\end{proof}

\begin{remark}[Size discipline]
In NBG, $\SQ$ and its operations are classes of set-coded tuples. In ZFC the same
notation can be used for definable class schemes. In a type theory, an ordinary
small type cannot contain every ordinal and hence every surreal. Work in a
fixed universe or a set-sized workspace and record the embeddings between
workspaces. None of the finite algebra requires a class of all classes, and
none of the analytic constructions requires class-indexed summation.
\end{remark}

\section{Division, norm, and vector geometry}
\label{sec:algebra}
Identify the pure imaginary subspace with $F^3$. For $q=a+u$ and $p=b+v$,
Hamilton multiplication becomes
\begin{equation}
 (a+u)(b+v)=ab-u\cdot v+av+bu+u\times v.
 \label{eq:dotcross}
\end{equation}
All dot and cross products here are finite algebraic expressions over $F$.

\begin{theorem}[Division and positive norm]
\label{thm:division}
Over any ordered field $F$,
\[
 \overline{pq}=\bar q\bar p,\qquad
 q\bar q=\bar q q=N(q),\qquad N(pq)=N(p)N(q).
\]
Moreover, $N(q)>0$ for $q\ne0$, so
\[
 q^{-1}=\frac{\bar q}{N(q)}.
\]
If $F$ is real closed, the $F$-valued modulus $\norm q=\sqrt{N(q)}$ satisfies
\[
 \norm{pq}=\norm p\norm q,\quad
 \norm{p+q}\le\norm p+\norm q,\quad
 \abs{\Rea(p\bar q)}\le\norm p\norm q.
\]
\end{theorem}
\begin{proof}
The conjugation identity follows from the relations on the basis. Then
$N(pq)=pq\bar q\bar p=pN(q)\bar p=N(p)N(q)$. A sum of squares in an ordered
field is zero only when all its summands are zero, proving division.
For Cauchy--Schwarz apply nonnegativity to
$N(p-sq)$, a quadratic polynomial in the central variable $s$. If $q\ne0$, its
minimum at $s=\Rea(p\bar q)/N(q)$ is nonnegative. Thus
$\Rea(p\bar q)^2\le N(p)N(q)$. Expanding $N(p+q)$ and taking the positive square
root gives the triangle inequality. No supremum, compactness, or completeness
is used.
\end{proof}

\begin{corollary}[Equality and inverse stability]
For nonzero $p,q$, equality in the triangle inequality holds exactly when
$q=cp$ for a positive central scalar $c$. For invertible $p,q$,
\begin{equation}
 p^{-1}-q^{-1}=p^{-1}(q-p)q^{-1},\qquad
 \norm{p^{-1}-q^{-1}}=
 \frac{\norm{p-q}}{\norm p\norm q}.
 \label{eq:inverse-lipschitz}
\end{equation}
\end{corollary}
\begin{proof}
Equality in the quadratic proof of Cauchy--Schwarz means central linear
dependence, and the positive sign in the real part selects $c>0$. The inverse
identity is verified by multiplication; its norm is exactly multiplicative.
\end{proof}

\begin{proposition}[Center, commutators, and ordering]
\label{prop:center}
The center of $\HH(F)$ is $F$. Furthermore,
\[
 [a+u,b+v]=2u\times v,\qquad
 \norm{[p,q]}\le2\norm{\Imag p}\norm{\Imag q}.
\]
There is no ordered-division-ring ordering on $\HH(F)$ extending the order of
$F$.
\end{proposition}
\begin{proof}
Equation \eqref{eq:dotcross} gives the commutator. A vector whose cross product
with each coordinate vector vanishes is zero, proving the assertion about the
center. The identity
$\norm{u\times v}^2=\norm u^2\norm v^2-(u\cdot v)^2$ proves the inequality.
In an ordered division ring every nonzero square is positive, contradicting
$\ii^2=-1$.
\end{proof}

\begin{example}[Several scales in one number]
For $q=\omega+\omega^{-1}\ii+\jj+\omega^{-\omega}\kk$,
\[
 N(q)=\omega^2+1+\omega^{-2}+\omega^{-2\omega},\qquad
 q^{-1}=\frac{\omega-\omega^{-1}\ii-\jj-\omega^{-\omega}\kk}
 {\omega^2+1+\omega^{-2}+\omega^{-2\omega}}.
\]
The coordinates can have arbitrarily different orders of magnitude without
compromising associativity, division, or the norm identities.
\end{example}

\section{Complex slices, conjugacy, and rotations}
\label{sec:slices}
Define the sphere of imaginary units
\[
 \Ss=\{I\in\HH(F):\Rea I=0,\ N(I)=1\}.
\]
Equivalently, $I^2=-1$. For each $I\in\Ss$ the slice $C_I=F+FI$ is a copy of
$F[\sqrt{-1}]$.

\begin{theorem}[Slice decomposition and centralizers]
\label{thm:slices}
Suppose $F$ is real closed. Every noncentral $q$ has a unique expression
\[
 q=a+rI,\qquad a\in F,\ r>0,\ I\in\Ss.
\]
Its centralizer is exactly $C_I$. Two different slices, unless their imaginary
units differ only by sign, intersect in $F$. Every noncentral element obeys
\[
 q^2-2a q+(a^2+r^2)=0,
\]
and this quadratic is its minimal polynomial over $F$.
\end{theorem}
\begin{proof}
Take $a=\Rea q$, $r=\norm{\Imag q}$, and $I=\Imag q/r$. The commutator formula
shows that an imaginary vector commutes with $I$ exactly when it is a scalar
multiple of $I$. This proves both centralizer statements. The quadratic follows
from $I^2=-1$, and no degree-one central polynomial can annihilate a noncentral
quaternion.
\end{proof}

\begin{proposition}[Conjugacy classes]
\label{prop:conjugacy}
Two quaternions are conjugate by an invertible quaternion if and only if they
have the same scalar part and norm square. The class of $a+rI$, $r>0$, is
\[
 [a+rI]=a+r\Ss.
\]
The conjugating quaternion may be chosen to have norm one.
\end{proposition}
\begin{proof}
Conjugation preserves the trace $2\Rea q$ and norm square. To send $I$ to $J$,
when $J\ne-I$ use $h=1-JI$: direct multiplication gives $hI=Jh$ and
$h\ne0$. When $J=-I$, choose a unit pure quaternion perpendicular to $I$; it
anticommutes with $I$ and supplies $h$. Normalize $h$ by its positive norm.
\end{proof}

\begin{proposition}[Matrix realization]
\label{prop:matrix}
Write $q=z+w\jj$ with $z,w\in F[\ii]$, and let the bar on this slice denote its
ordinary quadratic conjugation. Then
\begin{equation}
 \rho(q)=\begin{pmatrix}z&w\\-\bar w&\bar z\end{pmatrix}
 \label{eq:rho}
\end{equation}
is a faithful algebra homomorphism into $M_2(F[\ii])$, with
\[
 \rho(\bar q)=\rho(q)^*,\quad \det\rho(q)=N(q),\quad
 \operatorname{tr}\rho(q)=2\Rea q.
\]
The norm-one group is identified with $\SU(2,F[\ii])$. Scalar extension to
$F[\ii]$ splits the algebra:
$\HH(F)\otimes_F F[\ii]\simeq M_2(F[\ii])$.
\end{proposition}
\begin{proof}
Use $\jj z=\bar z\jj$ to obtain
$(z+w\jj)(u+v\jj)=(zu-w\bar v)+(zv+w\bar u)\jj$.
This is precisely matrix multiplication in \eqref{eq:rho}. The determinant and
star formulas follow directly. A unitary $2\times2$ matrix of determinant one
has the displayed form, by comparing its adjugate with its conjugate transpose.
Finally, the four images of the basis elements are linearly independent over
$F[\ii]$ and span its four-dimensional matrix algebra.
\end{proof}

\begin{theorem}[Spin description of rotations]
\label{thm:spin}
For real closed $F$, conjugation on pure imaginary vectors gives a surjective
homomorphism
\[
 \Sp(1,F):=\{u:N(u)=1\}\longrightarrow\SO(3,F),\qquad
 u\longmapsto(v\mapsto uv\bar u),
\]
with kernel $\{1,-1\}$. For $u=c+sI$, $c^2+s^2=1$,
\begin{equation}
 uv\bar u=(c^2-s^2)v+2s^2(I\cdot v)I+2cs(I\times v).
 \label{eq:Rodrigues}
\end{equation}
\end{theorem}
\begin{proof}
Expansion proves \eqref{eq:Rodrigues}; conjugation preserves dot products and
cross products, so its determinant is one. Its kernel consists of central unit
quaternions by \cref{prop:center}.

For surjectivity, an orthogonal $3\times3$ matrix of determinant one over a real
closed field has an eigenvector with eigenvalue one. Indeed, its cubic
characteristic polynomial has a real root; every real orthogonal eigenvalue is
$1$ or $-1$. If only a $-1$ root were initially found, its perpendicular plane
would have determinant $-1$, hence eigenvalues $1,-1$ (its characteristic
polynomial has discriminant $\operatorname{tr}^2+4>0$).
Normalize an eigenvector $I$. On its oriented perpendicular plane, the matrix
is $\bigl(\begin{smallmatrix}C&-S\\S&C\end{smallmatrix}\bigr)$ with
$C^2+S^2=1$. If $C\ne-1$, take
$c=\sqrt{(1+C)/2}$, $s=S/(2c)$; otherwise take $c=0,s=1$.
Equation \eqref{eq:Rodrigues} gives the required lift.
\end{proof}

Thus arbitrary-scale spatial rotations need no trigonometric value at an
infinite argument. They can be constructed entirely by rational operations and
positive square roots. A rational chart on the unit group is
\begin{equation}
 u=\frac{1-\norm x^2+2x}{1+\norm x^2},\qquad x\in F^3;
 \quad x=\frac{\Imag u}{1+\Rea u}\quad(u\ne-1).
 \label{eq:stereographic}
\end{equation}
The denominators are positive. This is stereographic parametrization, not a
claim that $\Sp(1,F)$ is topologically a real three-sphere.

\begin{corollary}[Automorphisms fixing the center]
Every $F$-algebra automorphism of $\HH(F)$ is inner, and
$\operatorname{Aut}_F\HH(F)\simeq\SO(3,F)$.
\end{corollary}
\begin{proof}
An automorphism takes $\ii,\jj,\kk$ to imaginary unit vectors with the same dot
and cross relations. They form an oriented orthonormal basis. Conversely such
a basis preserves Hamilton multiplication. Apply \cref{thm:spin}.
\end{proof}

\section{Polynomial equations and a constructive fundamental theorem}
\label{sec:polynomials}
A polynomial with quaternion coefficients must come with an evaluation
convention. We use a \emph{central indeterminate} $X$ and coefficients on the
right:
\[
 P(X)=\sum_{n=0}^d X^n a_n\in\HH(F)[X],\qquad
 P(q)=\sum_{n=0}^d q^n a_n.
\]
Multiplication in this polynomial ring will also be written $P*Q$ when it is
important to distinguish it from pointwise multiplication of evaluated
functions. Centrality of $X$ does not make the evaluation map multiplicative.
These conventions agree with the polynomial part of slice regularity
\cite{gp}.

\subsection{The product and division formulas}
\begin{proposition}[Twisted evaluation]
\label{prop:twisted}
When $P(q)\ne0$,
\begin{equation}
 (P*Q)(q)=P(q)\,Q\bigl(P(q)^{-1}qP(q)\bigr).
 \label{eq:twisted}
\end{equation}
When $P(q)=0$, one has $(P*Q)(q)=0$. Moreover, $P(q)=0$ if and only if
$P=(X-q)*R$ for some polynomial $R$.
\end{proposition}
\begin{proof}
Writing $Q=\sum X^n b_n$ and regrouping finite sums gives
$(P*Q)(q)=\sum_n q^nP(q)b_n$. If $P(q)\ne0$, then
$q^nP(q)=P(q)(P(q)^{-1}qP(q))^n$, proving \eqref{eq:twisted}. The zero case is
immediate. Synthetic division on the left by the monic polynomial $X-q$ gives
$P=(X-q)*R+c$ with constant remainder $c$. Evaluation of the product on the
right is zero at $q$, so $c=P(q)$.
\end{proof}

\subsection{Norm polynomials and conjugacy spheres}
Define $P^c(X)=\sum X^n\bar a_n$ and $P^s=P*P^c$. Write
\[
 P=P_0+P_1\ii+P_2\jj+P_3\kk,\qquad P_j\in F[X].
\]
Since the component polynomials commute,
\begin{equation}
 P^s=P^c*P=P_0^2+P_1^2+P_2^2+P_3^2\in F[X].
 \label{eq:symmetrization}
\end{equation}
Its degree is $2d$ and its leading coefficient is $N(a_d)>0$.

\begin{lemma}[Remainder on a sphere]
\label{lem:sphere-remainder}
Fix $a\in F$ and $b>0$, and put
$\Delta(X)=X^2-2aX+a^2+b^2$. Divide centrally:
\[
 P=\Delta Q+XA+B.
\]
On the sphere $a+b\Ss$, $P(q)=qA+B$. If $\Delta\mid P^s$, precisely one of the
following occurs: $A=B=0$, and the entire sphere is a zero set; or $A\ne0$, and
its unique zero is
\begin{equation}
 q=-BA^{-1},\qquad \Rea q=a,\quad N(q)=a^2+b^2.
 \label{eq:sphere-root}
\end{equation}
Conversely, a zero anywhere on this sphere implies $\Delta\mid P^s$.
\end{lemma}
\begin{proof}
The characteristic quadratic vanishes on the sphere. Conjugate and multiply
the remainder, then reduce modulo $\Delta$. The result is
\begin{equation}
 \bigl(2aN(A)+2\Rea(A\bar B)\bigr)X
       +N(B)-(a^2+b^2)N(A).
 \label{eq:remainder-norm}
\end{equation}
If $\Delta\mid P^s$, both central coefficients vanish. If $A=0$, positivity
forces $B=0$. Otherwise $q=-B\bar A/N(A)$ has precisely the scalar part and norm
in \eqref{eq:sphere-root}, so it belongs to the sphere and is the unique
solution of $qA+B=0$.

Conversely, if $qA+B=0$, then either $A=B=0$, or $B=-qA$. In the latter case
$N(B)=(a^2+b^2)N(A)$ and
$\Rea(A\bar B)=-aN(A)$. Equation \eqref{eq:remainder-norm} is zero.
\end{proof}

\begin{theorem}[Fundamental theorem for one-sided polynomials]
\label{thm:fta}
Over any real closed field $F$, every positive-degree $P\in\HH(F)[X]$ has a
zero in $\HH(F)$ and factors as
\[
 P=(X-q_1)*(X-q_2)*\cdots*(X-q_d)a_d.
\]
Its zero set meets at most $d$ conjugacy classes. In each noncentral class it is
either one point or the whole sphere.
\end{theorem}
\begin{proof}
The nonconstant central polynomial $P^s$ has a linear or irreducible quadratic
factor over $F$. A linear factor $X-r$ gives
$0=P^s(r)=N(P(r))$, so $P(r)=0$. Every irreducible quadratic over a real closed
field has, after making it monic, the form
$X^2-2aX+a^2+b^2$ with $b>0$. Apply \cref{lem:sphere-remainder} to obtain a zero.
Left division reduces the degree and yields the stated factorization by
induction.

Every noncentral zero class contributes its distinct quadratic factor to
$P^s$. Every central zero $r$ makes all four $P_j(r)$ vanish, so
$(X-r)^2\mid P^s$. Different classes supply relatively prime factors, each of
degree two. Since $\deg P^s=2d$, at most $d$ classes occur. The remainder formula
also proves the point-or-sphere alternative.
\end{proof}
The theorem is a real-closed-field version of established quaternionic
polynomial phenomena, with a proof not relying on a topological degree argument.
It supplies an actual finite root-recovery procedure from the central norm
polynomial. It does not promise that factoring central polynomials over an
arbitrary surreal presentation is computationally decidable.

\begin{example}[A factor does not give a pointwise root]
\label{ex:camshaft}
Let $P=(X-\ii)*(X-\jj)=X^2-X(\ii+\jj)+\kk$. Then
\[
 P(\ii)=0,\qquad P(\jj)=2\kk,\qquad P^s=(X^2+1)^2.
\]
The entire zero set is the single point $\ii$: modulo $X^2+1$ the nonzero linear
remainder has just one root. In contrast, $X^2+1$ vanishes on all of $\Ss$.
Degree bounds concern conjugacy classes or suitable multiplicities, not a
naive number of individual quaternion roots.
\end{example}

\begin{warning}[Not every noncommutative polynomial has a root]
The equation $q+\ii q\ii+1=0$ has no solution. For
$q=a+b\ii+c\jj+d\kk$ its left side is $1+2c\jj+2d\kk$. Thus
\cref{thm:fta} cannot be restated for arbitrary polynomials with coefficients
interspersed among all occurrences of the variable. The surquaternions are a
noncommutative division algebra, not an algebraically closed field.
\end{warning}

\subsection{Square roots}
\begin{proposition}[Explicit square roots]
\label{prop:sqrt}
Let $q=a+v\ne0$ and $R=\norm q$. If $q$ is not a negative central scalar, its
square root with positive scalar part is
\begin{equation}
 s=\sqrt{\frac{R+a}{2}}+
       \frac{v}{2\sqrt{(R+a)/2}}.
 \label{eq:sqrt}
\end{equation}
For a positive central scalar the same formula has $v=0$. For $q=-r^2<0$, all
square roots are $rI$, $I\in\Ss$. For noncentral $q$, there are exactly the two
roots $s,-s$; zero has only the root zero.
\end{proposition}
\begin{proof}
If $v\ne0$, then $R>\abs a$. Setting $x=\sqrt{(R+a)/2}$ gives
$\norm{v/(2x)}^2=(R-a)/2$, so squaring \eqref{eq:sqrt} gives $q$.
Every square root commutes with its square and therefore belongs to the same
slice as a noncentral $q$. In that field the two roots are exactly $s,-s$.
If a square has negative central value, its scalar-vector cross term forces its
scalar part to vanish, leaving precisely the stated sphere. Positivity treats
the remaining cases.
\end{proof}

\section{Finite-dimensional spectral theory}
\label{sec:spectral}
Let $V=\HH(F)^n$ be a right vector space with Hermitian product
$\langle x,y\rangle=\sum\bar x_jy_j$. Its squared norm is
$\langle x,x\rangle\in F_{\ge0}$. A matrix acts on the left. The adjoint $A^*$ is
the conjugate transpose. These conventions matter whenever eigenvalues or
coefficients are moved across vectors.

\begin{theorem}[Hermitian spectral theorem over a real closed field]
\label{thm:spectral}
If $A=A^*\in M_n(\HH(F))$ and $F$ is real closed, there are a quaternionic unitary
matrix $U$ and $\lambda_1,\ldots,\lambda_n\in F$ such that
\[
 U^*AU=\operatorname{diag}(\lambda_1,\ldots,\lambda_n).
\]
The matrix is positive definite exactly when every $\lambda_j>0$.
\end{theorem}
\begin{proof}
Extend \eqref{eq:rho} in blocks. The resulting $2n\times2n$ matrix $\rho(A)$ is
Hermitian over the algebraically closed field $F[\ii]$. Its characteristic
polynomial has a root $\lambda$ and a nonzero eigenvector $w$. Since
$w^*w>0$ and $w^*\rho(A)w$ is fixed by conjugation,
$\lambda=(w^*\rho(A)w)/(w^*w)$ belongs to $F$.

Every column $w\in F[\ii]^{2n}$ is obtained by stacking the first columns of
$\rho(x_j)$ for a unique $x\in\HH(F)^n$. Because $\lambda$ is central, the
matrix eigen-equation gives $Ax=x\lambda$. Normalize $x$ using the positive
square root of its norm square. Its quaternionic orthogonal complement is
$A$-invariant: if $\langle x,y\rangle=0$, then
$\langle x,Ay\rangle=\langle Ax,y\rangle=\lambda\langle x,y\rangle=0$.
Quaternionic Gram--Schmidt constructs an orthonormal basis of this complement,
and induction completes the diagonalization. The positivity criterion follows
by evaluating the diagonal quadratic form, including on basis vectors.
\end{proof}

\begin{corollary}[Polar decomposition and singular values]
An invertible $A\in M_n(\HH(F))$ has a unique decomposition
$A=UP$ with $U^*U=1$ and $P$ positive definite Hermitian. One has
$P=(A^*A)^{1/2}$. More generally every rectangular matrix has a singular value
decomposition with nonnegative central singular values.
\end{corollary}
\begin{proof}
Diagonalize $A^*A$ and take the positive square roots of its positive
eigenvalues, defining $P$. Then $U=AP^{-1}$ is unitary. A positive square root
of $A^*A$ commutes with that matrix. On each of its eigenspaces a positive
Hermitian square root can have only the positive eigenvalue $\sqrt\lambda$,
proving uniqueness. For a rectangular matrix, diagonalize $A^*A$; the images of
the eigenvectors with positive eigenvalues, divided by the corresponding
singular values, are orthonormal. Extend them to an orthonormal basis of the
codomain.
\end{proof}

\begin{example}[An infinitesimal avoided crossing]
For central $a,\delta\in F$ and $z\in\HH(F)$, the Hermitian matrix
\[
 A=\begin{pmatrix}a+\delta&z\\\bar z&a-\delta\end{pmatrix}
\]
satisfies $(A-aI)^2=(\delta^2+N(z))I$. Its eigenvalues are
$a\pm\sqrt{\delta^2+N(z)}$, with the coincident case included. Taking
$a=\omega$, $\delta=\omega^{-2}$, and $z=\omega^{-3}\jj$ gives the exact gap
$2\omega^{-2}\sqrt{1+\omega^{-2}}$. The infinite common offset and the much
smaller mixing scale cause no ambiguity in the algebraic spectrum.
\end{example}
No infinite-dimensional Hilbert-space spectral theorem is being asserted.
The result is finite linear algebra. More generally, fixed-dimensional
semialgebraic assertions can be transported through the theory of real closed
fields, subject to first-order expressibility \cite{swan}; assertions about all
infinite sequences or arbitrary analytic functions cannot be transported this
way.

\section{Normal forms, valuation, and residue quaternions}
\label{sec:valuation}
\subsection{Quaternion-valued Hahn series}
Writing surreal normal forms with $t^\gamma=\omega^{-\gamma}$ changes the
reverse-well-order convention to a well-order convention:
\[
 a=\sum_{\gamma\in S}a_\gamma t^\gamma,\qquad
 S\subseteq\No\text{ a well-ordered set},\quad a_\gamma\in\R.
\]
Combine the four supports of a surquaternion and put
$h_\gamma=a_\gamma+b_\gamma\ii+c_\gamma\jj+d_\gamma\kk\in\HH$.

\begin{theorem}[Normal-form and localization theorem]
\label{thm:normal}
Every surquaternion has a unique expression
\begin{equation}
 q=\sum_{\gamma\in S}t^\gamma h_\gamma,
 \qquad h_\gamma\in\HH\setminus\{0\},
 \label{eq:normal}
\end{equation}
with a set-sized well-ordered support. Addition is coefficientwise;
multiplication is ordered quaternion convolution:
\[
 (pq)_\eta=\sum_{\alpha+\beta=\eta}p_\alpha q_\beta.
\]
Each displayed coefficient sum is finite. Every set of surquaternions is
contained in one $D_\Gamma=\HH(\R((t^\Gamma)))$, where $\Gamma$ is a set-sized
divisible subgroup of $(\No,+)$. Conversely the coefficientwise identification
\[
 \HH(\R((t^\Gamma)))\simeq\HH((t^\Gamma))
\]
is an isomorphism of algebras with involution.
\end{theorem}
\begin{proof}
A finite union of well-ordered subsets of a linearly ordered set is
well-ordered. Uniqueness follows in each scalar coordinate. The ordinary Hahn
sumset lemma makes the convolution locally finite; the centrality of monomials
leaves the product of coefficients in its original order. For a set of
quaternions, union all its supports and take their $\Q$-linear span in $\No$.
This is a set-sized divisible ordered group $\Gamma$. Normal forms identify
its Hahn field with a subfield of $\No$; the standard real-closedness theorem
for real Hahn coefficients and divisible value group makes $K_\Gamma$ real
closed \cite{alling,ek}. The four-coordinate construction proves the final
isomorphism.
\end{proof}

\begin{remark}
All series in
\eqref{eq:normal} are normal-form/Hahn series. A workspace need not be closed
under Gonshor's global exponential or under a chosen scalar derivation; such
closure is an additional hypothesis when those operations are used.
\end{remark}

For $q\ne0$ let $\val(q)=\min\supp(q)$ and
$\operatorname{lc}(q)=h_{\val(q)}$. Put $\val(0)=+\infty$.

\begin{theorem}[Exact valuation and leading coefficients]
\label{thm:val}
For nonzero $p,q$,
\[
 \val(pq)=\val(p)+\val(q),\qquad
 \operatorname{lc}(pq)=\operatorname{lc}(p)\operatorname{lc}(q).
\]
Also
\[
 \val(p+q)\ge\min(\val(p),\val(q)),\quad
 \val(\bar q)=\val(q),\quad
 \val(N(q))=2\val(q),\quad \val(\norm q)=\val(q).
\]
In coordinates, $\val(a+b\ii+c\jj+d\kk)=\min(\val(a),\val(b),\val(c),\val(d))$.
\end{theorem}
\begin{proof}
The smallest exponent in a product is the sum of the two smallest exponents,
and its quaternion coefficient is the product of two nonzero elements of
$\HH$. This cannot vanish. In $q\bar q$ the coefficient is
$N(\operatorname{lc}(q))>0$, proving the remaining multiplicative assertions.
The additive and coordinate assertions follow directly from supports.
\end{proof}
The valuation is an additive order-of-magnitude measure, not the norm. For
example, $1$ and $2+\ii$ have the same valuation but different norms.

\subsection{Finite elements and standard part}
Let
\[
 \OO_D=\{q:\val(q)\ge0\},\qquad \mm_D=\{q:\val(q)>0\}.
\]
These are exactly the finite-norm and infinitesimal-norm quaternions,
respectively, where ``finite'' means bounded by some positive real number.
Coefficient extraction at exponent zero gives
\begin{equation}
 \st:\OO_D\longrightarrow\HH,
 \qquad \OO_D/\mm_D\simeq\HH.
 \label{eq:residue}
\end{equation}
This is a surjective homomorphism of algebras with involution. Every finite
$q$ decomposes uniquely as $q=\st(q)+\varepsilon$ with $\varepsilon\in\mm_D$.
The residue is the ordinary quaternion division algebra, not $\C$ and not a
commutative field.

\begin{proposition}[Multiplicative normal form]
\label{prop:mult-normal}
Every nonzero $q$ has a unique expression
\begin{equation}
 q=t^\gamma h(1+\varepsilon),\qquad
 \gamma=\val(q),\quad h=\operatorname{lc}(q)\in\HH^\times,
 \quad\varepsilon\in\mm_D.
 \label{eq:mult-normal}
\end{equation}
Moreover,
\[
 q^{-1}=\left(\sum_{n\ge0}(-\varepsilon)^n\right)h^{-1}t^{-\gamma}.
\]
The sum is a Hahn sum. The finite unit group is a split extension of
$\HH^\times$ by $1+\mm_D$, with conjugation of the latter by constant
quaternions. Monomial factors are central.
\end{proposition}
\begin{proof}
Remove the leading monomial and multiply on the left by $h^{-1}$. The remainder
has strictly positive valuation. The support lemma in the next section
justifies the geometric series, whose product with $1+\varepsilon$ is one by
formal cancellation. Reduction modulo $\mm_D$ supplies the split group map.
For example, the product of $h(1+\varepsilon)$ and $k(1+\delta)$ is
$hk(1+k^{-1}\varepsilon k)(1+\delta)$, displaying the action explicitly.
\end{proof}

\section{Summability and the topology that must not be confused with it}
\label{sec:support}
\begin{definition}[Hahn-summable family]
A set-indexed family $(q_j)_{j\in J}$ is Hahn-summable if the union of its
supports is well-ordered and each exponent belongs to the supports of only
finitely many $q_j$. Its sum is obtained by adding those finitely many
quaternion coefficients at each exponent. All uses of an infinite sum below
refer to this definition unless an ordinary real or complex integral is
explicitly named.
\end{definition}

\begin{lemma}[Positive-support lemma]
\label{lem:neumann}
Let $S$ be a well-ordered subset of the positive cone of an ordered abelian
group. Its generated additive monoid
\[
 S^*=\{s_1+\cdots+s_n:n\ge0,\ s_j\in S\}
\]
is well-ordered. For a fixed exponent $\gamma$, only finitely many finite
ordered words $(s_1,\ldots,s_n)$ over $S$ have sum $\gamma$, counting all lengths
at once. Consequently, every noncommutative formal power series in finitely
many variables with real quaternion coefficients can be evaluated at finitely
many elements of $\mm_D$, retaining the order of all factors.
\end{lemma}
\begin{proof}
The standard Neumann support lemma gives these assertions; here is the reason
its all-lengths conclusion is available. A well-ordered alphabet is
well-quasi-ordered. Higman's finite-word lemma says that its finite words are
well-quasi-ordered under subword embedding with entrywise domination. Such an
embedding can only increase the sum of a word. It increases it strictly unless
no entries are inserted and no entries increase, because all entries are
positive. A strictly decreasing sequence of word sums would contradict the
well-quasi-order property. Distinct words with one fixed sum form an antichain,
so there are only finitely many. This proves both conclusions (with the usual
choice principle used to characterize a non-well-order by a decreasing
sequence).

For evaluation, take the finite union of the supports of the arguments. A
contribution at exponent $\gamma$ is specified by a word over this union, a
choice of one of the finitely many variables at each position, and its
coefficient term. There are only finitely many such contributions. The union
of possible exponents lies in $S^*$.
\end{proof}
The repository's \path{Surreal/HahnSeries/Neumann.lean} proves this same
all-lengths form \cite{repo}. It is important
that the proof works for arbitrary ordered abelian groups, not only rank-one
or Archimedean groups.

\begin{warning}[Positive valuation is a summability certificate, not cofinality]
If $\val(\varepsilon)=\delta>0$, then $\val(\varepsilon^n)=n\delta$.
The exponents $n\delta$ need not be cofinal in the value group. In the
lexicographically ordered group $\Z\times\Z$, the sequence $n(0,1)$ never exceeds
$(1,0)$. Nevertheless the geometric series is Hahn-summable. Thus a valid Hahn
identity need not be the limit of its partial sums in the valuation topology of
the workspace.
\end{warning}

\subsection{The full surreal fine topology}
Use balls $B(q,r)=\{p:\norm{p-q}<r\}$ with arbitrary positive surreal radii to
define the fine topology on $\SQ$. Here this means the ball-neighbourhood
scheme; no collection whose elements are all open classes is required.

\begin{theorem}[Set discreteness]
\label{thm:fine}
Every set of surquaternions is closed and discrete in the full fine topology.
Every convergent set-indexed net in that topology is eventually constant.
Every continuous map from an ordinary connected real interval into $\SQ$ is
constant.
\end{theorem}
\begin{proof}
Given a set $A$ and a point $q$, the nonzero distances
$\{\norm{a-q}:a\in A,\ a\ne q\}$ form a set of positive surreals. Every set of
positive surreals has a strictly positive lower bound: take a surreal larger
than all their reciprocals and invert it. A smaller ball about $q$ therefore
misses $A\setminus\{q\}$. This proves discreteness for $q\in A$ and closedness
for $q\notin A$.

A set-indexed net and its proposed limit have a set of values; isolate the
limit in that set to obtain eventual equality. A map from a real interval has
set-sized image, which is discrete in its subspace topology. A continuous image
of a connected space in a discrete space is a singleton.
\end{proof}
In particular, no sequence of distinct real rational approximations converges
finely to $\sqrt2$, even though it converges in $\R$. Nor do the partial sums of
$\sum\omega^{-n}$ converge finely to their Hahn sum.

\begin{remark}[A workspace has two different topologies]
The intrinsic order or valuation topology on $K_\Gamma$ uses radii or values
available in that workspace. Its topology as a subspace of the full fine
surreal topology is discrete, because $K_\Gamma$ is a set. These need not agree.
For example, in $\R((t^\Q))$ the intrinsic valuation admits ordinary geometric
partial-sum convergence, whereas the same set is discrete inside $\No$ with
all surreal radii. Even intrinsically, a higher-rank group can have the
noncofinality phenomenon just described.
\end{remark}
A different, deliberately coarse topology on finite surquaternions is the
standard-part topology, pulled back from $\HH\simeq\R^4$ by \eqref{eq:residue}.
It is not Hausdorff: points with the same standard part cannot be separated.
Coefficientwise analysis over ordinary domains uses that ordinary domain as
its geometric base. It does not pretend that its paths are fine-continuous
surreal paths.

\section{Infinitesimal exponential, logarithm, and rotation groups}
\label{sec:lie}
\subsection{An exact noncommutative local calculus}
For $x,y\in\mm_D$ define
\begin{equation}
 \exph(x)=\sum_{n\ge0}\frac{x^n}{n!},\qquad
 \logh(1+y)=\sum_{n\ge1}\frac{(-1)^{n+1}y^n}{n}.
 \label{eq:local-exp-log}
\end{equation}
The subscript emphasizes Hahn evaluation rather than a limit of partial sums.
Each individual series lies in the commutative subalgebra generated by its one
argument and the central scalars.

\begin{theorem}[Local exp--log equivalence]
\label{thm:local-exp}
The maps in \eqref{eq:local-exp-log} are inverse bijections
\[
 \exph:\mm_D\longleftrightarrow1+\mm_D:\logh.
\]
They commute with quaternion conjugation and inner automorphisms. For $x\ne0$,
\[
 \val(\exph(x)-1)=\val(x),\qquad
 \operatorname{lc}(\exph(x)-1)=\operatorname{lc}(x).
\]
If $x$ and $y$ commute, $\exph(x+y)=\exph(x)\exph(y)$.
\end{theorem}
\begin{proof}
The positive-support lemma evaluates the ordinary formal identities
$\log(\exp X)=X$ and $\exp(\log(1+Y))=1+Y$ coefficientwise. The coefficient of
the leading exponent comes only from the linear term. Conjugation reverses
products, but it does not alter powers of a single element. Inner automorphisms
preserve products and valuation. Multiplying a summable family by the fixed
factors $h$ and $h^{-1}$ preserves summability by the Hahn sumset lemma; the
individual supports may change. When $x,y$ commute, the binomial identity
proves the exponential law; the common support certificate permits the double
sum to be regrouped.
\end{proof}
For noncommuting arguments, the correct group law on $\mm_D$ is
\[
 \BCH(x,y)=\logh\bigl(\exph(x)\exph(y)\bigr).
\]
The usual formal identity yields
\begin{align}
 \BCH(x,y)={}&x+y+\tfrac12[x,y]
 +\tfrac1{12}[x,[x,y]]+\tfrac1{12}[y,[y,x]]\notag\\
 &-\tfrac1{24}[y,[x,[x,y]]]+\text{terms of degree at least five}.
 \label{eq:BCH}
\end{align}
Here ``degree'' counts occurrences of $x$ and $y$, not powers of an assumed
single numerical scale. If both valuations are at least $\delta>0$, the last
remainder has valuation at least $5\delta$; a common support certificate remains
necessary for the whole infinite expression. Formula \eqref{eq:BCH} can be
obtained directly by multiplying formal exponentials and expanding the formal
logarithm. This derivation, rather than any analytic BCH convergence theorem,
is the one used here.

\begin{proposition}[Adjoint series]
For $x\in\mm_D$ and any fixed $q\in D_\Gamma$,
\begin{equation}
 \exph(x)q\exph(-x)=
 \sum_{n\ge0}\frac{\ad_x^n(q)}{n!},\qquad \ad_x(q)=[x,q].
 \label{eq:adjoint}
\end{equation}
The series is Hahn-summable, including when $q$ is infinite.
\end{proposition}
\begin{proof}
The formal identity follows by collecting terms in
$\sum_{r,s}x^rq(-x)^s/(r!s!)$. Its supports lie in
$\supp(q)+S^*$, where $S=\supp(x)$; the Hahn sumset lemma combined with
\cref{lem:neumann} gives local finiteness. Noncommutativity is retained in the
words $x^rqx^s$.
\end{proof}

\subsection{The group of infinitesimal rotations}
Let $U=\{u\in\OO_D:N(u)=1\}$ and
$U_1=U\cap(1+\mm_D)$. Reduction gives an exact split sequence
\[
 1\longrightarrow U_1\longrightarrow U
 \xrightarrow{\st}\Sp(1,\R)\longrightarrow1.
\]
Every unit quaternion is finite, so this describes the whole norm-one group
in a workspace, and the corresponding class statement holds over $\No$.

\begin{theorem}[Logarithms of infinitesimal rotations]
\label{thm:rotation-log}
Hahn exponential and logarithm restrict to mutually inverse maps
\[
 \exph:\Imag(\mm_D)\longleftrightarrow U_1:\logh.
\]
On the left the group law is BCH, not addition. If $u=u_0w$ with
$u_0=\st(u)\in\Sp(1,\R)$, then $w\in U_1$ is unique. Thus every finite-scale
unit rotation is an ordinary rotation together with a unique infinitesimal
BCH correction, once this multiplication order is fixed.
\end{theorem}
\begin{proof}
For pure imaginary $x$, $\overline{\exph(x)}=\exph(-x)$, so the norm is one.
If $u\in U_1$, then $\bar u=u^{-1}$, and
$\overline{\logh u}=\logh(\bar u)=\logh(u^{-1})=-\logh u$.
Hence its logarithm is pure imaginary. The factorization is
$w=u_0^{-1}u$, which has standard part one and norm one.
\end{proof}

\begin{theorem}[Associated graded rotation algebra]
\label{thm:graded}
For $\gamma\in\Gamma_{>0}$ set
$U_{\ge\gamma}=\{u\in U_1:\val(u-1)\ge\gamma\}$ and define
$U_{>\gamma}$ similarly. Leading coefficients give
\[
 U_{\ge\gamma}/U_{>\gamma}\simeq(\Imag\HH,+).
\]
If $u=\exph(x)$, $v=\exph(y)$ with $x,y$ pure infinitesimals of respective
valuations $\alpha,\beta$ and leading coefficients $X,Y$, then
\[
 uvu^{-1}v^{-1}=1+t^{\alpha+\beta}[X,Y]
     +\text{terms of valuation greater than }\alpha+\beta.
\]
In particular its valuation above one is exactly $\alpha+\beta$ when
$X\times Y\ne0$, and the induced graded bracket is $[X,Y]=2X\times Y$.
\end{theorem}
\begin{proof}
The product of $1+a$ and $1+b$ has leading correction $a+b$, since the product
$a b$ has larger valuation. The norm-one condition forces the leading
correction to be pure imaginary. Every such coefficient is realized by
$\exph(t^\gamma X)$. Expanding the group commutator, its first nonzero formal
term is $[x,y]$; each subsequent term contains both arguments and at least one
additional positive-valuation argument. This proves the formula, including
possible cancellation when $[X,Y]=0$.
\end{proof}

\begin{example}[Rotations at unequal scales]
Let $x=t^\alpha\ii$ and $y=t^\beta\jj$, with $\alpha,\beta>0$. Then
\[
 \BCH(x,y)=t^\alpha\ii+t^\beta\jj+t^{\alpha+\beta}\kk
 -\tfrac13t^{2\alpha+\beta}\jj
 -\tfrac13t^{\alpha+2\beta}\ii+\cdots.
\]
The leading failure of commutativity appears at the sum of the two scales.
There is no requirement that $\alpha$ and $\beta$ be rationally comparable.
For a vector $v$, the leading infinitesimal action is
$\exph(x)v\exph(-x)=v+2t^\alpha(\ii\times v)+\cdots$.
\end{example}

\section{Finite angles and polar coordinates at arbitrary scale}
\label{sec:angles}
Let $\OO_{\No}$ and $\mm_{\No}$ denote the finite and infinitesimal surreal
scalars. If $\theta=r+\eta$ with $r\in\R$ and $\eta\in\mm_{\No}$, define
\begin{align*}
 \cos\theta&=\cos r\sum_{n\ge0}\frac{(-1)^n\eta^{2n}}{(2n)!}
 -\sin r\sum_{n\ge0}\frac{(-1)^n\eta^{2n+1}}{(2n+1)!},\\
 \sin\theta&=\sin r\sum_{n\ge0}\frac{(-1)^n\eta^{2n}}{(2n)!}
 +\cos r\sum_{n\ge0}\frac{(-1)^n\eta^{2n+1}}{(2n+1)!}.
\end{align*}
Real sine and cosine are evaluated first. Their real Taylor sums are not
misrepresented as Hahn-summable families of constants. These finite-angle
functions agree with the restricted analytic surreal extensions used in
\cite{ek} and in the repository's trigonometric report \cite{repo}.

\begin{proposition}[Finite phase theorem]
\label{prop:finitephase}
For a fixed $I^2=-1$ the map
\[
 \theta\longmapsto\cos\theta+I\sin\theta
\]
is a surjective homomorphism from $(\OO_{\No},+)$ onto the unit circle in
$\No+\No I$, with kernel $2\pi\Z$.
\end{proposition}
\begin{proof}
The addition and norm identities are formal Taylor identities with a common
positive support. If a phase is one, reducing to its standard part gives
$r\in2\pi\Z$; the infinitesimal exponential is injective, forcing $\eta=0$.

Conversely write a unit-circle point as $z=c+Is$, with finite $c,s$. Its
standard part in this abstract slice is an ordinary circle point
$z_0=\cos r+I\sin r$. Then $z_0^{-1}z$ is one plus an infinitesimal. Its logarithm
is pure imaginary in the same slice, hence $I\eta$ for a scalar infinitesimal
$\eta$. Thus $z=\cos(r+\eta)+I\sin(r+\eta)$.
The argument takes standard parts of the two slice coordinates, not of the
symbol $I$; it therefore also works for a surreal imaginary unit.
\end{proof}

\begin{theorem}[Polar representation]
\label{thm:polar}
Every nonzero surquaternion is expressible as
\[
 q=\rho(\cos\theta+I\sin\theta),\qquad
 \rho=\norm q>0,\quad \theta\in\OO_{\No},\quad I^2=-1.
\]
For noncentral $q$, one can uniquely choose
$I=\Imag q/\norm{\Imag q}$ and $0<\theta<\pi$ with
$\cos\theta=\Rea q/\rho$ and $\sin\theta=\norm{\Imag q}/\rho$.
Positive central elements use $\theta=0$; negative ones use $\theta=\pi$ with
arbitrary $I$.
\end{theorem}
\begin{proof}
The normalized scalar and imaginary norm lie on the unit circle with the
second coordinate positive. Apply \cref{prop:finitephase}. The sign of sine
selects the open upper semicircle and the unique representative in $(0,\pi)$.
This includes the cases where the angle or its defect from $\pi$ is
infinitesimal, as is seen directly from its Taylor expansion.
\end{proof}
An infinite magnitude $\rho$ thus does not require an infinite angle. The
surreal scale resides in the radius; even an imaginary direction with
infinitesimal coordinate components is represented by a finite angle. The
usual rotation angle for the spin action is $2\theta$, as follows from
\eqref{eq:Rodrigues} and the finite-angle double-angle identities.

\section{A global radial exponential and its logarithm fibres}
\label{sec:global}
The local exponential is forced by its formal series. A global phase on all
surreal scalars requires additional structure. We now make that structure
explicit instead of treating $\sum q^n/n!$ as an everywhere-defined Hahn sum.

\subsection{The Ehrlich--Kaplan normalization}
Split the scalar normal form into its purely infinite and finite parts:
\[
 x=x_{\mathrm{PI}}+\fin(x),\qquad
 \fin(x)=r+\eta\in\OO_{\No}.
\]
The projection $\fin$ is additive, but is not multiplicative. The purely
infinite part has only positive powers of $\omega$. In particular,
$\fin(\omega)=0$. Put
\begin{equation}
 \Sin x=\sin(\fin x),\qquad \Cos x=\cos(\fin x).
 \label{eq:global-phase}
\end{equation}
These are the canonical scalar phase functions associated with the initial
integer part
\[
 \Oz=\No_{\mathrm{PI}}+\Z.
\]
Indeed, $2\pi\Oz=\No_{\mathrm{PI}}+2\pi\Z$, so reducing modulo this integer-part
period group amounts exactly to discarding the purely infinite part and then
reducing the finite angle. This is the normalization of Ehrlich--Kaplan,
Section 11 \cite{ek}. It gives $\Sin\omega=0$, $\Cos\omega=1$ and, on each
complex slice, a surjective exponential with period group $2\pi I\Oz$.
These are consequences of a specified normalization, not consequences of
ordered-field algebra alone.

Let $\exp_{\No}:\No\to\No_{>0}$ be the Gonshor exponential and $\log_{\No}$ its
inverse \cite{gonshor,ek}. It extends the real exponential and obeys the usual
additive-to-multiplicative law.

\begin{definition}[Canonical radial surquaternion exponential]
For $q=a+v$, $r=\norm v$, define
\begin{equation}
 \ExpC(q)=
 \begin{cases}
 \exp_{\No}(a)\bigl(\Cos r+(v/r)\Sin r\bigr),&r>0,\\
 \exp_{\No}(a),&r=0.
 \end{cases}
 \label{eq:global-exp}
\end{equation}
\end{definition}
The construction is radial: it uses the actual norm $r$ and actual direction
$v/r$ of the imaginary part, then applies the chosen scalar phase to $r$.
Discarding the purely infinite parts of the three vector coordinates first
would define a different, generally non-equivariant operation.

\begin{theorem}[Basic global exponential properties]
\label{thm:globalexp}
The map \eqref{eq:global-exp} is a surjection $\SQ\to\SQ^\times$. It satisfies
\[
 \norm{\ExpC(q)}=\exp_{\No}(\Rea q),\qquad
 \ExpC(-q)=\ExpC(q)^{-1},\qquad
 \ExpC(\bar q)=\overline{\ExpC(q)}.
\]
For $h\ne0$,
\[
 \ExpC(hqh^{-1})=h\ExpC(q)h^{-1}.
\]
It restricts on every slice $\No+\No I$ to the canonical surcomplex
exponential, and satisfies $\ExpC(p+q)=\ExpC(p)\ExpC(q)$ whenever $p,q$ commute.
On $\mm_D$ it agrees with $\exph$.
\end{theorem}
\begin{proof}
The norm, inverse, and conjugation formulas follow from the scalar phase
identities. Inner conjugation preserves $a,r$ and carries the direction $I$ to
$hIh^{-1}$, proving equivariance. On a slice, writing $q=a+Ib$ shows that
\eqref{eq:global-exp} equals $\exp_{\No}(a)(\Cos b+I\Sin b)$, using the oddness of
$\Sin$ when $b<0$. Commuting noncentral quaternions lie in one slice, by
\cref{thm:slices}; central cases cause no change.

Given $q\ne0$, take its finite-angle polar representation from
\cref{thm:polar}. Then
$\ExpC(\log_{\No}\norm q+I\theta)=q$, proving surjectivity. For an infinitesimal
$q=a+v$, the central scalar $a$ commutes with $v$, and the even and odd powers
of $v$ are exactly the Hahn cosine and sine series at $r$. Expanding the
commuting product gives the local series.
\end{proof}

\begin{theorem}[Complete logarithm fibres]
\label{thm:logfibres}
For noncentral $q=\rho(\cos\theta+I\sin\theta)$ in the canonical polar form
$0<\theta<\pi$, all its preimages are exactly
\begin{equation}
 \ExpC^{-1}(q)=
 \{\log_{\No}\rho+I(\theta+2\pi z):z\in\Oz\}.
 \label{eq:logfibre-nonreal}
\end{equation}
For $\rho>0$ the central fibres are
\begin{align}
 \ExpC^{-1}(\rho)&=
 \{\log_{\No}\rho+I(2\pi z):I^2=-1,\ z\in\Oz\},\label{eq:positive-fibre}\\
 \ExpC^{-1}(-\rho)&=
 \{\log_{\No}\rho+I\pi(2z+1):I^2=-1,\ z\in\Oz\}.
 \label{eq:negative-fibre}
\end{align}
There is no preimage of zero.
\end{theorem}
\begin{proof}
The norm determines the scalar part of every preimage. The image of a
quaternion lies in its own slice. If the image is noncentral, its centralizer
is the unique slice of that image, so the preimage is
$\log_{\No}\rho+Ib$ for a scalar $b$. The finite phase theorem together with
\eqref{eq:global-phase} makes its phase fibre exactly
$b\in\theta+2\pi\Oz$. When the image is central, its imaginary part vanishes,
so its radial phase is $1$ or $-1$; every axis is allowed. These give
\eqref{eq:positive-fibre} and \eqref{eq:negative-fibre}, including repetitions
of the scalar logarithm when $z=0$ in the positive case. Norms are always
strictly positive, excluding zero.
\end{proof}
The fibre above one should not be called the kernel of a global group
homomorphism: the additive quaternion group is abelian and the multiplicative
group is not. The exponential law is asserted only for commuting arguments.
Nor do the fibre descriptions settle the robustness questions for alternative
surcomplex exponentials posed in \cite{ek}.

\begin{example}[Infinite radius is not coordinatewise phase reduction]
With $q=\omega\ii+\pi\jj$,
\[
 r=\sqrt{\omega^2+\pi^2}
   =\omega+\frac{\pi^2}{2\omega}+\bigO(\omega^{-3}),
\]
so
\[
 \ExpC(q)=1+\frac{\pi^2}{2\omega}\ii+\bigO(\omega^{-2}).
\]
In contrast, $\ExpC(\omega\ii)\ExpC(\pi\jj)=-1$. Here and below, order symbols
for Hahn expansions bound the valuation of the remainder; they are not
sequence limits. The example also shows why the direction and radius must be
formed before the purely infinite phase is discarded.
\end{example}

\section{Differential of the radial exponential}
\label{sec:differential}
Differentiation of a quaternionic function with respect to four scalar
coordinates is not the same operation as a scalar-field derivation. We first
use the ordinary $\No$-linear Fr\'echet definition, expressed by inequalities:
$L$ is the derivative of $f$ at $q$ when, for every surreal $\epsilon>0$, some
surreal $\delta>0$ satisfies
\[
 0<\norm h<\delta\quad\Longrightarrow\quad
 \norm{f(q+h)-f(q)-L(h)}<\epsilon\norm h.
\]
This definition is meaningful as a class scheme even though set-indexed fine
limits are degenerate. It does not assert a real Banach-space structure.

\begin{theorem}[Radial derivative and critical spheres]
\label{thm:exp-derivative}
Let $q=a+rI$ with $r>0$. Write a variation as
$h=s+\alpha I+w$, where $s,\alpha\in\No$ and $w$ is pure imaginary and
perpendicular to $I$. Put $E=\exp_{\No}(a)$, $C=\Cos r$, $S=\Sin r$. Then
\begin{equation}
 D\ExpC(q)[h]=E\left(s(C+IS)+\alpha(-S+IC)+\frac{S}{r}w\right).
 \label{eq:exp-derivative}
\end{equation}
At a central point $a$, the derivative is $h\mapsto\exp_{\No}(a)h$. As a
four-dimensional $\No$-linear map,
\begin{equation}
 \det D\ExpC(a+rI)=\exp_{\No}(4a)\left(\frac{\Sin r}{r}\right)^2
 \quad(r>0).
 \label{eq:exp-jacobian}
\end{equation}
It is singular precisely for positive $r\in\pi\Oz$, and has rank two there.
\end{theorem}
\begin{proof}
Fix $q$. For a sufficiently small variation $u=\alpha I+w$ relative to $r$,
positive-square-root expansion gives
\[
 \norm{rI+u}=r+\alpha+\bigO(\norm u^2/r),\qquad
 \frac{rI+u}{\norm{rI+u}}
     =I+\frac{w}{r}+\bigO(\norm u^2/r^2).
\]
These are Taylor identities on infinitesimal relative increments, justified
by the positive-support lemma; alternatively the norm expansion follows by
rationalizing the difference of two square roots. The constants in their
norm bounds may depend on the fixed surreal $r$.

For infinitesimal scalar $\eta$, the additive phase law gives
$\Cos(r+\eta)=C-S\eta+\bigO(\eta^2)$ and
$\Sin(r+\eta)=S+C\eta+\bigO(\eta^2)$, even when $r$ is infinite.
Likewise $\exp_{\No}(a+s)=E(1+s+\bigO(s^2))$. Combining the three expansions
proves \eqref{eq:exp-derivative} with a remainder bounded by
$M_q\norm h^2$ on a sufficiently small ball, for some positive surreal $M_q$.
Shrinking the ball further to radius below $\epsilon/(M_q+1)$ proves the
stated fine derivative. At $r=0$, use the Hahn exponential at the small
increment $h$ and the commuting central scalar $a$.

On the scalar--radial plane the matrix is
$E\bigl(\begin{smallmatrix}C&-S\\S&C\end{smallmatrix}\bigr)$, of determinant
$E^2$. On the two transverse coordinates it is $(ES/r)$ times the identity.
This proves \eqref{eq:exp-jacobian}. Finally, $\Sin r=0$ precisely when
$\fin r\in\pi\Z$, equivalently $r\in\pi\Oz$. The scalar--radial block remains
invertible, so a singular derivative has rank exactly two.
\end{proof}

\begin{example}[A critical sphere of infinite radius]
At $q=\omega I$, the normalized phase has $C=1,S=0$. Thus
\[
 D\ExpC(\omega I)[s+\alpha I+w]=s+\alpha I.
\]
Every transverse first-order variation vanishes, while the scalar and radial
ones survive. For a central formal small parameter $\eta$ and $J\perp I$,
\[
 \ExpC(\omega I+\eta J)
    =1+\frac{\eta^2}{2\omega}I
       +\frac{\eta^3}{2\omega^2}J+\text{higher terms}.
\]
This is not obtained by setting the exponential of each imaginary coordinate
independently. It is a radial critical-point phenomenon at an infinite scale.
\end{example}
The determinant formula extends the familiar finite-radius criticality of the
quaternion exponential, but the allowed radii now include a proper class of
infinite integer-part periods. A nonsingular fine derivative by itself is not
invoked as an unrestricted inverse-function theorem; the support-based theorem
in \cref{sec:lifting} is the positive existence tool used in this article.

\section{Slice and Fueter analysis with explicit coefficient spaces}
\label{sec:analysis}
There is more than one useful quaternionic analogue of holomorphy. The two we
use have different differential operators and different product rules. Their
classical relationship is developed in \cite{gp}. Our additional requirement
is that every infinite operation be defined in a specified coefficient space.

\subsection{Slice power-series germs}
Let $a_n\in\HH$ and let $q\in\mm_D$. Define
\[
 f(q)=\sum_{n\ge0}q^n a_n.
\]
This is Hahn-summable regardless of the ordinary real radius of convergence of
the coefficients: infinitesimal support, rather than a real analytic estimate,
is doing the work. More generally we may take $a_n\in\OO_D$ if all their
supports lie in one well-ordered set $S\subseteq\Gamma_{\ge0}$. The supports of
the evaluated terms then lie in $S+T^*$, where $T=\supp(q)$, and the support
lemmas give local finiteness. Merely knowing $\val(a_n)\ge0$ individually is
not a substitute for this common-support hypothesis.

\begin{proposition}[Slice representation and derivative]
\label{prop:representation}
For an admissible series $f$ as above, infinitesimal central $x,y$, and any
$I,J\in\Ss$ in the workspace, one has
\begin{equation}
 f(x+yJ)=\tfrac12(1-JI)f(x+yI)+\tfrac12(1+JI)f(x-yI).
 \label{eq:representation}
\end{equation}
On a fixed slice, coefficientwise formal differentiation satisfies
\[
 (\partial_x+I\partial_y)f(x+yI)=0,\qquad
 \partial_s f(q)=\sum_{n\ge1}nq^{n-1}a_n.
\]
The formal convolution product has the twisted evaluation law
\eqref{eq:twisted} whenever the nonzero value used for conjugation is
invertible; when it is zero, the product value is zero.
\end{proposition}
\begin{proof}
For central $x,y$, the powers of $x+yI$ have the form $A_n(x,y)+I B_n(x,y)$
with real integer-coefficient polynomials $A_n,B_n$, independent of $I$.
Summing gives $f(x+yI)=A+IB$. Solving the two equations with $I$ and $-I$
produces \eqref{eq:representation}. The Cauchy--Riemann and derivative identities
hold for every monomial and then by common-support regrouping.

For the product formula, repeat the proof of \cref{prop:twisted}. Inner
conjugation leaves the valuation and norm of the infinitesimal argument
unchanged. A direct support certificate also follows from the fixed factors
$f(q)^{\pm1}$ and the supports of the original evaluated terms, so the
regrouping has the same validity as the finite formula.
\end{proof}
A nonzero formal germ with \emph{ordinary quaternion coefficients} has a first
nonzero coefficient $a_m$. Then, for every nonzero infinitesimal $q$,
$f(q)=q^m(a_m+\varepsilon)$ with $\varepsilon\in\mm_D$, so it is nonzero.
Hahn-valued coefficients allow genuinely displaced infinitesimal zeros, such
as the zero of $q-t^\gamma\ii$; thus these two coefficient classes must not be
silently interchanged.

\subsection{Common-domain coherent functions}
For analytic continuation and contours, fix an ordinary circular quaternionic
domain $\Omega\subseteq\HH$ and a well-ordered set $S\subseteq\Gamma$. Consider
formal expressions
\begin{equation}
 f=\sum_{\gamma\in S}t^\gamma f_\gamma,
 \qquad f_\gamma:\Omega\to\HH
 \text{ ordinary stem-induced slice-regular functions}.
 \label{eq:coherent}
\end{equation}
The domain $\Omega$ is common to all coefficients. Here slice regularity is understood in the stem-function sense, on a
conjugation-invariant complex domain:
$f_\gamma(x+Iy)=A_\gamma(x,y)+I B_\gamma(x,y)$, where the quaternion-valued
components satisfy the Cauchy--Riemann equations and the even/odd parity that
makes this definition independent of the sign of $I$.

For finite $x,y$ with standard part in that domain, evaluate $A_\gamma,B_\gamma$
by their ordinary Taylor germs at the standard point. Substitute the positive
Hahn corrections, and then sum over $\gamma$. The support certificate is
$S+T^*$, where $T$ is the union of the supports of the finitely many scalar
increments. The additional factor $I$ is a fixed finite quaternion and adds
only its well-ordered support. This defines the corresponding function at
surreal directions as well. Uniqueness of ordinary Taylor germs ensures that
shrinking a neighbourhood does not change the value.

The slice product in this space is induced by multiplying stem functions. It
is not generally the pointwise product of quaternion values. Coefficientwise
differentiation and this product preserve the stated function space, with
support sets transformed by finite sums and sumsets. A family of analytic
germs with no common ordinary domain would be a different space and would not
support the contour statement that follows.

\begin{theorem}[Coefficientwise Cauchy formula on a fixed slice]
\label{thm:cauchy}
Fix an ordinary imaginary unit $I$ and an ordinary positively oriented
piecewise smooth Jordan contour $C$ whose interior and boundary lie in
$\Omega\cap C_I$. Each $f_\gamma$ is therefore holomorphic on a neighbourhood
of that closed interior. If $q_0$ is an ordinary interior point
and $q=q_0+\eta$ with infinitesimal $\eta\in K_\Gamma+IK_\Gamma$, then
\begin{equation}
 f(q)=\frac{1}{2\pi}\int_C (s-q)^{-1}\,ds_I\,f(s),\qquad ds_I=-I\,ds.
 \label{eq:cauchy}
\end{equation}
The integral and its infinitesimal kernel expansion are defined
coefficientwise. In particular, \eqref{eq:cauchy} is not an integral along a
fine-continuous surquaternion path.
\end{theorem}
\begin{proof}
On $C_I$, split an ordinary quaternion-valued holomorphic function into two
complex-valued holomorphic components, with a fixed perpendicular unit on the
right. The ordinary complex Cauchy formula proves \eqref{eq:cauchy} for each
$f_\gamma$ at $q_0$ and gives all its derivatives there. Expand
\[
 (s-q_0-\eta)^{-1}=\sum_{n\ge0}\eta^n(s-q_0)^{-n-1}.
\]
The contour stays a positive ordinary distance from $q_0$. Each coefficient
integral is therefore an ordinary, well-defined integral and equals the
corresponding Taylor coefficient. Supports of the total family lie in
$S+\supp(\eta)^*$ and are locally finite. Thus coefficientwise integration
reconstructs exactly the defined value $f(q)$.
\end{proof}
No claim is made that testing an arbitrary surquaternion function at ordinary
points determines it. Such uniqueness is a property of the specified coherent
coefficient representation, not of all functions on $\SQ$.

\subsection{Fueter regularity and the Laplacian bridge}
For $q=a+b\ii+c\jj+d\kk$, define the left Cauchy--Fueter operator and the
coordinate Laplacian by
\[
 \mathcal D=\partial_a+\ii\partial_b+\jj\partial_c+\kk\partial_d,
 \qquad \Delta_4=\partial_a^2+\partial_b^2+\partial_c^2+\partial_d^2.
\]
A function is left Fueter regular when $\mathcal Df=0$. The identity function
$q\mapsto q$ is slice regular but has $\mathcal Dq=-2$, so the theories are
not interchangeable.

\begin{theorem}[Polynomial and coherent Fueter transfer]
\label{thm:fueter}
For every right-coefficient slice polynomial $P$ over a real closed field,
\[
 \mathcal D\Delta_4 P=0.
\]
The same identity holds coefficientwise for coherent slice-regular functions
of the form \eqref{eq:coherent}.
\end{theorem}
\begin{proof}
Away from the real axis, write $q=a+rI$ and $f=A(a,r)+IB(a,r)$, with
$A_a=B_r$ and $A_r=-B_a$. Direct differentiation in the three imaginary
coordinates gives
\[
 \mathcal Df=(A_a-B_r-2B/r)+I(B_a+A_r),
\]
and, because $A$ and $B$ are harmonic in $(a,r)$,
\[
 \Delta_4 f=C+ID,\qquad
 C=2A_r/r,\quad D=2B_r/r-2B/r^2.
\]
Now
\[
 C_a-D_r-2D/r=0,\qquad D_a+C_r=0
\]
by the two Cauchy--Riemann equations and their derivatives. Hence
$\mathcal D\Delta_4 f=0$ for $r>0$. For a polynomial the identity extends over
the real axis as a polynomial identity; it can also be read as an identity in
the four scalar coordinates and consequently holds over every real closed
field. For an ordinary slice-regular function the coefficients are real
analytic, so the same identity extends across real points by continuity and
analyticity. Apply this separately to every $f_\gamma$ in \eqref{eq:coherent}.
All differentiations in that space are coefficientwise by definition.
\end{proof}
This gives a meaningful bridge between two quaternionic analytic theories
without claiming a nonexistent global fine-topological Cauchy theory. Other
classical results may admit similar coefficientwise versions, but their
function space, common domain, and boundary hypotheses have to be stated anew.

\section{Support-controlled implicit lifting and square-root conditioning}
\label{sec:lifting}
This section gives a mechanism for solving perturbed quaternionic equations.
Its essential hypothesis is nonsingularity of the \emph{full real-coordinate
Jacobian}. A nonzero derivative on a single complex slice is not enough.

\subsection{A finite system with a nonsingular real reduction}
For a scalar vector $x=(x_1,\ldots,x_d)$, set
$\val(x)=\min_j\val(x_j)$, with $\val(0)=+\infty$. Write $\OO_K$ and
$\mm_K$ for the valuation ring and its maximal ideal in $\KG$.

\begin{theorem}[Support-controlled implicit lifting]
\label{thm:lifting}
Let $P\in\OO_K[X_1,\ldots,X_d]^d$, and let $x_0\in\R^d$ satisfy
\[
 \st P(x_0)=0,\qquad
 J_0=\left(\frac{\partial\st P_i}{\partial X_j}(x_0)\right)_{i,j}
 \in\GL_d(\R).
\]
There is exactly one $h\in\mm_K^d$ such that $P(x_0+h)=0$.
Let $S\subset\Gamma_{>0}$ be the union of the supports of the positive-order
parts of the finitely many coefficients of $P(x_0+X)$. Then
\[
 \supp(h_j)\subseteq S^*\setminus\{0\}\qquad(1\leq j\leq d),
\]
where $S^*$ is the additive monoid generated by $S$. If $P(x_0)\ne0$, then
\[
 \val(h)=\val(P(x_0)).
\]
No countable convergence or cofinality assumption on $\Gamma$ is required.
\end{theorem}
\begin{proof}
Write the finitely many coefficient perturbations as $e_1,\ldots,e_m\in\mm_K$.
Replace them by independent central variables $Y_1,\ldots,Y_m$. This gives a
polynomial system $\mathcal P(X,Y)$ over $\R$, with
$\mathcal P(0,0)=0$ and $X$-Jacobian $J_0$ at $(0,0)$.

There is a unique formal solution $H(Y)\in(Y)\R[[Y]]^d$. Here is the complete
recursion needed for that assertion. If all homogeneous terms of $H$ below
degree $n$ have been determined, the degree-$n$ part of
$\mathcal P(H(Y),Y)$ is $J_0H_n+R_n$, where $R_n$ depends only on the already
determined terms. Set $H_n=-J_0^{-1}R_n$. Each degree involves only finitely
many monomials, and this determines a formal solution and proves its formal
uniqueness.

Evaluate $Y_j=e_j$. By \cref{lem:neumann}, the family of all evaluated
monomials is summable: its supports lie in $S^*$, and each exponent has only
finitely many contributing words. Thus $h=H(e_1,\ldots,e_m)$ exists as a Hahn
vector, has the asserted support, and satisfies $P(x_0+h)=0$ because finite
polynomial operations respect this evaluation.

To prove uniqueness among \emph{all} infinitesimal vectors, not merely among
the solutions obtained by this substitution, let $h,h'\in\mm_K^d$. Polynomial
divided differences give
\[
 P(x_0+h)-P(x_0+h')=(J_0+E)(h-h'),\qquad E\in M_d(\mm_K).
\]
The determinant of $J_0+E$ has nonzero real constant coefficient, so the
matrix is invertible. Consequently two roots coincide.

Finally, put $r=P(x_0)$. The equation for $h$ is
$r+J_0h+E_1h+R_{\geq2}(h)=0$, with $E_1$ of positive order and all
coefficients of $R_{\geq2}$ integral. If $h\ne0$, both error terms have
valuation strictly greater than $\val(h)$. An invertible real matrix preserves
vector valuation, so cancellation at the least exponent forces
$\val(r)=\val(h)$. If $r=0$, uniqueness instead gives $h=0$.
\end{proof}

The argument is a support-explicit form of nonsingular Hensel lifting. Its
extra usefulness here is that quaternion multiplication is polynomial in four
scalar coordinates. Any finite system of quaternion equations, including
expressions with coefficients inserted on both sides of the unknowns, becomes
such a scalar system. Thus \cref{thm:lifting} applies as soon as its
real-coordinate Jacobian is nonsingular. This does \emph{not} give a root theorem
for every noncommutative polynomial expression; the counterexample in
\cref{sec:polynomials} still applies.

\subsection{The exact Sylvester operator for squaring}
Let $q_0=a+rI\ne0$, with $r>0$, $I^2=-1$. Every increment has the unique
orthogonal decomposition
\[
 h=h_{\parallel}+h_{\perp},\qquad
 h_{\parallel}=\frac{h-IhI}{2}\in F+FI,\qquad
 h_{\perp}=\frac{h+IhI}{2},
\]
where $h_{\perp}$ is purely imaginary and perpendicular to $I$.

\begin{proposition}[Square-root linearization and condition numbers]
\label{prop:sylvester}
The derivative of $q\mapsto q^2$ at $q_0$ is
\begin{equation}
 L_{q_0}(h)=q_0h+hq_0
          =2q_0h_{\parallel}+2a h_{\perp}.
 \label{eq:sylvester}
\end{equation}
Its determinant as an $F$-linear map on $F^4$ is
\[
 \det L_{q_0}=16a^2N(q_0).
\]
For $a\ne0$ it has inverse
\begin{equation}
 L_{q_0}^{-1}(w)=(2q_0)^{-1}w_{\parallel}+(2a)^{-1}w_{\perp}.
 \label{eq:sylvester-inverse}
\end{equation}
Its singular values are $2\norm{q_0}$ twice and $2|a|$ twice. If $a=0$, its
kernel is exactly the two-dimensional perpendicular plane.
\end{proposition}
\begin{proof}
The parallel component commutes with $I$ and the perpendicular component
anticommutes with it. This proves the displayed formula by expansion. On the
parallel two-plane, left multiplication by $2q_0$ is a similarity of scale
$2\norm{q_0}$ and determinant $4N(q_0)$. On the perpendicular two-plane the
map is multiplication by $2a$. The remaining assertions follow immediately.
\end{proof}

A purely imaginary square root can have a perfectly nonzero slice derivative
$2q_0$ while the four-dimensional derivative has a two-dimensional kernel.
This is the linear-algebraic explanation for the spheres of roots of a
negative scalar. It is also the reason that one cannot import a complex
simple-root test without modification.

\begin{example}[A noncommuting displaced square root]
Let $t$ be a positive infinitesimal monomial and $q_0=1+\ii$. There is a unique
root close to $q_0$ of
\[
 q^2=q_0^2+t\jj.
\]
Its first terms are
\begin{equation}
 q=1+\ii+\frac{t}{2}\jj
       +\frac{t^2}{16}(1-\ii)-\frac{t^3}{32}\jj+\bigO(t^4).
 \label{eq:sqrt-example}
\end{equation}
Here $\bigO(t^4)$ denotes a Hahn remainder of valuation at least $4\val(t)$,
not an ordinary limiting process. Substituting $q=q_0+\sum_{n\geq1}t^nh_n$
gives the exact recursion
\[
 L_{q_0}(h_1)=\jj,\qquad
 L_{q_0}(h_n)=-\sum_{j=1}^{n-1}h_jh_{n-j}\quad(n\geq2).
\]
The support lies in $\N\val(t)$, and \cref{thm:lifting} justifies the whole
series rather than only the displayed truncation.
\end{example}

\subsection{A quantified near-singular lifting region}
The following result keeps track of the small perpendicular singular value.
It is useful even when $q_0$ itself has no nonsingular real reduction.

\begin{theorem}[A sufficient conditioned square-root threshold]
\label{thm:conditioned}
Let $a\in\mm_K$ be nonzero, put $\kappa=\val(a)>0$ and $q_0=a+\ii$.
For $e\in\DG$ with $\sigma=\val(e)>2\kappa$, there is exactly one solution
of
\[
 q^2=q_0^2+e
\]
with $\val(q-q_0)>\kappa$. Writing $h=q-q_0$, one has
\begin{equation}
 \val(h)\geq\sigma-\kappa,\qquad
 \val\bigl(h-L_{q_0}^{-1}e\bigr)\geq2\sigma-3\kappa.
 \label{eq:conditioned-bounds}
\end{equation}
The condition $\sigma>2\kappa$ is asserted as a sufficient condition, not
as an optimal threshold for every direction of $e$.
\end{theorem}
\begin{proof}
Write $h=az$. The equation $L_{q_0}h+h^2=e$ becomes
\[
 z+aL_{q_0}^{-1}(z^2)=a^{-1}L_{q_0}^{-1}e.
\]
The real-coordinate matrix of $aL_{q_0}^{-1}$ has integral entries: its
parallel block is multiplication by $a/(2q_0)$, and its perpendicular block
is $1/2$. Moreover, every entry of $L_{q_0}^{-1}$ has valuation at least
$-\kappa$. The right side therefore has positive valuation at least
$\sigma-2\kappa$. The system in $z$ has real reduction vanishing at $0$ with
Jacobian the identity. Apply \cref{thm:lifting} to get a unique infinitesimal
$z$ and the bound $\val(z)\geq\sigma-2\kappa$. This gives the first estimate
for $h$. The exact equation
$h-L_{q_0}^{-1}e=-L_{q_0}^{-1}(h^2)$ and the multiplicativity of valuation give
the second. The condition $\val(h)>\kappa$ is precisely $z\in\mm_D$, which
also proves the stated uniqueness region.
\end{proof}

The theorem identifies the relevant loss as that of the \emph{inverse
linearized operator}. The determinant alone records a product of four scales
and obscures which perturbation directions are ill-conditioned. Parallel
errors and perpendicular errors can behave very differently.

\section{Scalar derivations and differential quaternion algebra}
\label{sec:derivations}
A field derivation is a different object from the derivative of a function of
four scalar variables. This distinction becomes especially important over
$\No$: the derivation differentiates the \emph{surreal coefficients themselves}.

\subsection{Extension and classification of derivations}
Let $\delta:F\to F$ be a derivation, meaning an additive map satisfying
$\delta(xy)=x\delta(y)+\delta(x)y$.

\begin{theorem}[Coefficientwise extension and inner part]
\label{thm:derivations}
The formula
\[
 \delta_H(a+b\ii+c\jj+d\kk)
 =\delta(a)+\delta(b)\ii+\delta(c)\jj+\delta(d)\kk
\]
defines the unique derivation of $\HF$ extending $\delta$ and annihilating
$\ii,\jj,\kk$. It satisfies
\[
 \delta_H(\bar q)=\overline{\delta_H(q)},\qquad
 \delta N(q)=2\Rea(\delta_H(q)\bar q),
 \qquad
 \delta\norm{q}=\frac{\Rea(\delta_H(q)\bar q)}{\norm{q}}\quad(q\ne0).
\]
Every derivation $D$ of $\HF$ has the unique form
\begin{equation}
 D(q)=\delta_H(q)+[w,q],\qquad w\in\Imag\HF,
 \label{eq:derivation-decomp}
\end{equation}
where $\delta$ is its restriction to the center $F$.
\end{theorem}
\begin{proof}
Hamilton multiplication has integer structure constants, which proves the
Leibniz rule for the coefficientwise extension. Commutation with conjugation
is immediate. Differentiate $N(q)=q\bar q$ and then $\norm{q}^2=N(q)$ to get
the two scalar formulas.

A derivation preserves the center: differentiating $aq-qa=0$ shows that
$D(a)$ commutes with every $q$. Subtract the coefficientwise extension of this
central restriction. The result $E$ is $F$-linear. Differentiating
$\ii^2=\jj^2=\kk^2=-1$ shows that $E(\ii),E(\jj),E(\kk)$ are imaginary
and perpendicular to the corresponding basis vectors. Differentiating the
anticommutation relations shows that the $3\times3$ matrix of $E$ on the
imaginary subspace is skew-symmetric. Every such matrix is uniquely
$v\mapsto2w\times v$ for an imaginary $w$. By the commutator formula this
is $[w,v]$, and both maps vanish on scalars. Thus $E=\ad_w$. Uniqueness follows
because an imaginary element commuting with all of $\HF$ is zero.
\end{proof}

This is a concrete quaternionic form of the inner-derivation property of
central simple algebras. No infinite-dimensional algebra theorem is needed
for the proof.

\subsection{The Berarducci--Mantova choice on the surreal field}
Berarducci and Mantova construct a distinguished surreal derivation $\partial$
that is strongly additive, compatible with the surreal exponential, has
constant field $\R$, and is surjective; their normalization has
$\partial\omega=1$ \cite[Theorems A and B]{bm}. These are external scalar
results, not consequences of the quaternion construction.

\begin{corollary}
The coefficientwise extension $\partial_H$ to $\SQ$ is surjective and has
constant algebra exactly $\HH$. In particular, every surquaternion has a
coefficientwise antiderivative.
\end{corollary}
\begin{proof}
An element is annihilated exactly when each of its four coordinates lies in
$\ker\partial=\R$. Choose one scalar antiderivative for each of its four
coordinates to obtain surjectivity.
\end{proof}

A set-sized Hahn workspace need not be stable under $\partial$. Statements
about a differential workspace require closure under that derivation or an
explicit enlargement. Also, surjectivity of $\partial_H$ is not a theorem
that arbitrary nonlinear quaternionic differential equations have solutions.

For a strongly additive derivation and $x\in\mm_D$, when the scalar
workspace is stable under the relevant operations, termwise differentiation
gives
\begin{equation}
 \delta_H\exph(x)=
 \sum_{n\geq1}\frac1{n!}\sum_{j=0}^{n-1}
    x^j\delta_H(x)x^{n-1-j}.
 \label{eq:derivative-exp-noncomm}
\end{equation}
The family on the right is support-controlled: it is built from words in the
positive support of $x$ with one distinguished factor from
$\supp\delta_H(x)$. Neumann finiteness and the finite antidiagonals of sums
of well-ordered sets justify the displayed regrouping. In general it is not
$\exph(x)\delta_H(x)$. That simpler expression is valid when
$[x,\delta_H(x)]=0$.

\subsection{Differential rotations}
\begin{proposition}[Angular velocity identity]
Let $u\in\HF$ satisfy $N(u)=1$, and set
$\xi=\delta_H(u)u^{-1}$. Then $\xi$ is purely imaginary. For every $x\in\HF$,
\begin{equation}
 \delta_H(uxu^{-1})=[\xi,uxu^{-1}]+u\delta_H(x)u^{-1}.
 \label{eq:angular-velocity}
\end{equation}
In particular, for a constant imaginary $x$, the derivative of its rotated
value is $2\xi\times(uxu^{-1})$.
\end{proposition}
\begin{proof}
Differentiate $u\bar u=1$ to get $\xi+\bar\xi=0$. The inverse rule is
$\delta_H(u^{-1})=-u^{-1}\delta_H(u)u^{-1}$. Apply Leibniz' rule to
$uxu^{-1}$ and collect terms.
\end{proof}
This formula applies to surreal differential coefficients without pretending
that the parameter is an ordinary real-time path into the fine topology.
It provides an algebraic angular-velocity equation, not an existence theorem
for trajectories.

\section{Exact computation and a formalization architecture}
\label{sec:computation}
The four-coordinate representation is computationally inexpensive relative
to its scalar backend. The difficulty is not Hamilton multiplication; it is
what the backend can prove about surreal coefficients and infinite supports.

\subsection{Representations and certified operations}
A useful implementation represents a surquaternion by a quadruple of scalar
objects together with their common field or workspace certificate. Addition,
conjugation, multiplication, norm square, and inversion reduce to finitely
many scalar operations. A magnitude additionally needs a positive square root
in the ordered scalar field. The implementation must not replace this
$F$-valued magnitude by a machine floating-point norm.

For Hahn objects, maintain well-ordered support data, a leading exponent,
and an ordinary quaternion leading coefficient. Inversion first removes the
leading monomial and coefficient and then uses the positive-order geometric
series. For an exact infinite object, a support certificate and a coefficient
rule are part of its meaning. A finite truncation should state the valuation
of its omitted remainder. Such truncations need not provide all coefficients
below an arbitrary cutoff in finitely many steps; an infinite support can
have infinitely many exponents below that cutoff.

Polynomial root computation has a precise algebraic route. Form the central
norm polynomial $P^s$, factor it over the real closed scalar field into linear
and irreducible quadratic factors, and compute the two-term remainder modulo
each quadratic. This determines whether a root class is a point or a whole
sphere and recovers an isolated point by $-BA^{-1}$. Actual effectiveness still
depends on the scalar field's factorization facilities. This is a mathematical
algorithm relative to an exact scalar oracle, not an algorithm for arbitrary
surreal inputs supplied without finite descriptions.

For implicit lifting, the finite formal recursion in
\cref{thm:lifting} is implementable up to a chosen formal degree. If supports
have several independent scales, a formal-degree cutoff and a valuation
cutoff are different specifications. A correct program records both the
computed expression and the support or valuation information proving what
has been omitted.

\subsection{Reuse the quaternion type, not four unrelated scalar copies}
Mathlib already provides a generic quaternion algebra, conjugation, norm
square, and division-ring instances under ordered-field hypotheses
\cite{mathlib}. The initial formalization should instantiate and extend those
objects, rather than duplicate their multiplication proofs. In particular,
a commutative complexification abstraction alone does not encode the
noncommuting quaternion units.

The inspected repository already contains generic algebra and Hahn support
lemmas, but its README explicitly distinguishes those prerequisites from a
construction of the surreal field and the normal-form bridge \cite{repo}.
A surquaternion extension must preserve that distinction. The following are
\emph{proposed modules}, not files that this article claims to have added or
compiled.

\begin{center}
\begin{tabular}{@{}p{0.32\textwidth}p{0.61\textwidth}@{}}
\toprule
Proposed layer & Main statements and dependencies\\
\midrule
Quaternion algebra & Conjugacy, slices, one-sided polynomial reduction, finite
spectral theory; generic ordered or real closed scalar field.\\[4pt]
Quaternion Hahn series & Coefficientwise equivalence, multiplicative valuation,
residue algebra; four-coordinate transport of support lemmas.\\[4pt]
Infinitesimal Lie theory & Exponential, logarithm, BCH and unit-group filtration;
Neumann finiteness for all words.\\[4pt]
Support-controlled lifting & Formal multivariable recursion and its evaluation;
real-coordinate Jacobian certificate.\\[4pt]
Coherent slice analysis & Common ordinary domains, coefficientwise contour
operations, representation and Fueter identities.\\[4pt]
Full surreal bridge & Class or universe discipline, normal forms, scalar
exponential, finite part, and the chosen global phase.\\
\bottomrule
\end{tabular}
\end{center}

The first five layers can be meaningfully proved in set-sized generic
workspaces. The last layer is not obtained by declaring a small type whose
elements are all surreal numbers. In a dependent-type formalization, an
explicit universe convention or a family of compatible set-sized models is
necessary. The ordered magnitude also calls for care: the usual real-valued
normed-space interfaces are not interchangeable with an $F$-valued magnitude.

\subsection{What the accompanying verification does}
The archive includes a SymPy script checking finite polynomial, rational, and
truncated-series identities. Its inputs are exact indeterminates and rational
numbers, not approximate floating-point samples. The recorded output names
each check, so that its scope can be inspected. In particular, it tests
noncommutative product ordering, the square-root linearization, and truncated
Lie-series formulas that are easy places for sign errors.

These checks do not machine-check the proofs about proper classes,
well-ordered supports, real-closed-field root existence, or infinite families.
No Lean implementation or formal verification of this article is asserted.
The proposed architecture is intended to make those future proof obligations
visible rather than to conceal them behind a successful symbolic calculation.

\section{Limitations, research directions, and synthesis}
\label{sec:conclusion}
The algebraic construction is robust: surquaternions form a central division
algebra of dimension four over the surreal class field, in the class-schematic
sense made precise earlier. Its elementary identities do not depend on
analytic choices. More delicate constructions have more specific domains.

The polynomial root theorem concerns \emph{one-sided} polynomials with a
central formal variable, not arbitrary expressions with coefficients on both
sides. Finite spectral theorems concern finite-dimensional matrices, not a
Hilbert-space theory over a proper class. The local exponential is a
support-controlled operation on infinitesimals. The global radial exponential
uses the established scalar normalization and is not a homomorphism on the
whole additive quaternion algebra. Fine differentiation does not create
ordinary continuous paths or justify a general inverse-function theorem.
Coefficientwise integration is defined over ordinary contours and is not an
integral over a fine-continuous surreal contour.

Three directions appear especially substantive. First, quaternionic zero
classes can change from spheres to isolated points under small coefficient
perturbations. A deformation theory should retain the whole finite
conjugacy-class algebra, not only track a list of scalar roots. The explicit
remainder $XA+B$ and the singular Sylvester operator identify the data such a
theory must control. Second, repeated quaternionic eigenvalues call for a
support-aware perturbation theory of eigenspaces and spectral projections;
the finite Hermitian theorem supplies existence but not a uniform series for
an arbitrarily chosen eigenbasis through a collision. Third, a comparison of
global radial phase extensions should separate their shared infinitesimal
Lie theory from their infinite-radius critical sets. None of these directions
is claimed solved by the constructions above.

The main lesson is a constructive one. Almost all finite quaternion algebra
survives over a real closed field; surreal normal forms then add an exact
leading-term and residue calculus; positive support supplies a noncommutative
analytic calculus on infinitesimals; and common-domain coefficient families
permit a separate, ordinary-analytic layer. Keeping those layers explicit
produces a substantial theory without asking a single topology or a single
meaning of ``convergence'' to do incompatible jobs.

\appendix
\section{Notation and conventions}
\begin{longtable}{@{}p{0.24\textwidth}p{0.7\textwidth}@{}}
\toprule
Symbol & Meaning\\
\midrule
\endhead
$F$, $\No$ & A real closed field; the surreal proper class, respectively.\\
$\HF$, $\SQ$ & Hamilton quaternion algebra over $F$; its surreal instance.\\
$\HH$ & Ordinary quaternions over $\R$.\\
$\bar q$, $N(q)$, $\norm q$ & Quaternion conjugate, norm square $q\bar q$,
and positive square root of the norm square.\\
$\Rea q$, $\Imag q$ & Scalar part and pure-imaginary part, not maps to
ordinary real coefficients unless the input is ordinary.\\
$\Ss$ & Imaginary units $I$ with $I^2=-1$; a two-sphere over $F$.\\
$\KG$, $\DG$ & $\R((t^\Gamma))$ and $\HH((t^\Gamma))\cong\HH_{\KG}$,
for a divisible set-sized ordered subgroup $\Gamma\subseteq\No$.\\
$t^\gamma$ & The surreal monomial $\omega^{-\gamma}$; supports in the
$t$-convention are well ordered.\\
$\val$ & Least support exponent; $\val(0)=+\infty$.\\
$\OO_D$, $\mm_D$, $\st$ & Nonnegative-order subring, positive-order ideal,
and its ordinary quaternion residue map.\\
$S^*$ & All finite sums from a positive well-ordered support $S$, including $0$.\\
$P*Q$, $P^c$, $P^s$ & Central-variable polynomial product, coefficient
conjugation, and central symmetrization $P*P^c$.\\
$\exph$, $\logh$ & Hahn exponential on infinitesimals and logarithm on
principal units.\\
$\fin x$ & Real-plus-infinitesimal part of the scalar normal form.\\
$\Oz$ & Canonical surreal integer part: purely infinite part plus an ordinary
integer.\\
$\ExpC$ & Global radial quaternion exponential using the Ehrlich--Kaplan
scalar phase normalization.\\
$\mathcal D$, $\delta_H$ & Fueter differential operator in four variables;
coefficientwise extension of a field derivation.\\
\bottomrule
\end{longtable}

\section{Reproducibility and reading routes}
\label{app:reproduce}
The source is a standalone LaTeX article with an embedded bibliography. It
requires no external figures, bibliography database, or proprietary fonts.
From the archive directory, compile with
\begin{verbatim}
latexmk -pdf -interaction=nonstopmode -halt-on-error surquaternions.tex
\end{verbatim}
Run the finite identity checks with
\begin{verbatim}
python3 verify.py
\end{verbatim}
The script needs SymPy. It writes a JSON record giving the installed SymPy
version and the result of each check. In the delivered run, all 31 named
checks passed using SymPy 1.14.0 and Python 3.13.5. Its truncation order is an explicit
finite parameter, and no finite truncation is presented as a proof of an
infinite-series identity.

For an algebra-first reading, follow Sections 2--7, then Section 14. For
infinitesimal geometry, follow Sections 7--10 and then Section 14. For the
exponential at infinite arguments, read Sections 10--12 together: omitting
the scalar normalization in Section 11 changes the subject. For analysis and
formalization, read Sections 8, 13, 15, and 16 with their different notions of
differentiation kept separate.

The repository snapshot used for orientation is
\texttt{39f2be6667ade51bca2b45daa47e289d69c09764}; the inspected sources include
its root README, document catalogue, trigonometry report overview, and the
Hahn Neumann-support module. They establish the intended interfaces and the
scope of the repository's existing formalization. The present proofs are
written out independently rather than treating unrefereed repository reports
as a substitute for proof. Literature inspection was targeted, not an
exhaustive priority search. In particular, familiar real-closed-field,
quaternionic, and scalar-surreal constructions are not advertised as new.

\clearpage
\phantomsection
\addcontentsline{toc}{section}{References}
\begin{thebibliography}{99}
\raggedright

\bibitem{conway}
J. H. Conway, \emph{On Numbers and Games}, 2nd ed., A K Peters, 2001.

\bibitem{gonshor}
H. Gonshor, \emph{An Introduction to the Theory of Surreal Numbers},
London Mathematical Society Lecture Note Series 110, Cambridge University
Press, 1986.

\bibitem{alling}
N. L. Alling, \emph{Foundations of Analysis over Surreal Number Fields},
North-Holland, 1987.

\bibitem{ek}
P. Ehrlich and E. Kaplan, \emph{Surreal ordered exponential fields},
arXiv:2002.07739, 2020, especially Sections 2, 3, and 11.
\url{https://arxiv.org/abs/2002.07739}.

\bibitem{bm}
A. Berarducci and V. Mantova, \emph{Surreal numbers, derivations and
transseries}, arXiv:1503.00315, 2015, especially Theorems A and B.
\url{https://arxiv.org/abs/1503.00315}.

\bibitem{voight}
J. Voight, \emph{Quaternion Algebras}, Springer, 2021. Author-maintained
open-access text and corrections:
\url{https://jvoight.github.io/quat.html}.

\bibitem{gp}
R. Ghiloni and A. Perotti, Slice regular functions on real alternative algebras,
\emph{Advances in Mathematics} \textbf{226} (2011), 1662--1691.
\url{https://arxiv.org/abs/1008.4318}.

\bibitem{swan}
R. G. Swan, \emph{Tarski's Principle and the Elimination of Quantifiers},
expository notes.
\url{https://www.math.uchicago.edu/~swan/expo/Tarski.pdf}.

\bibitem{mathlib}
The mathlib community, \emph{Mathlib.Algebra.Quaternion}, online documentation,
consulted 21 September 2026.
\url{https://leanprover-community.github.io/mathlib4_docs/Mathlib/Algebra/Quaternion.html}.

\bibitem{repo}
V. Reshetnikov, \emph{Surreal}, GitHub repository, snapshot
\texttt{39f2be6667ade51bca2b45daa47e289d69c09764}, inspected 21 September 2026.
\url{https://github.com/VladimirReshetnikov/Surreal/tree/39f2be6667ade51bca2b45daa47e289d69c09764}.
Specific orientation sources: \texttt{README.md}, \texttt{docs/README.md},
\texttt{docs/surcomplex/trigonometry/README.md}, and
\texttt{Surreal/HahnSeries/Neumann.lean}.

\end{thebibliography}
\end{document}
